% ===========================================================================
%  PLATAFORMA UNOORIA — Respostas ao questionário de validação do TR
%  (Liga Digital, versão 1 do questionário, 30/10/2026) — PU-EST-003
%  Documento autonomo. Compilar a partir desta pasta:
%      tectonic -k respostas-questionario-tr.tex
% ===========================================================================
\documentclass[11pt, a4paper, oneside, openany]{report}

% ============================================================================
%  estilo.tex — identidade visual, tipografia e componentes do documento
%
%  Tipografia: Roboto (texto e titulos) + Roboto Mono (codigo). As duas vem do
%  pacote `roboto`, entao o documento compila com pdfLaTeX, XeLaTeX ou Tectonic
%  sem fontspec explicito.
%
%  Paleta: os tokens `--unooria-*` do @unooria/design-system, os MESMOS valores
%  que `packages/ui/styles/plataforma.css` usa como fallback. Cor nova neste
%  documento so entra se existir como token no pacote.
% ============================================================================

% --- Codificacao e idioma ---------------------------------------------------
\usepackage[T1]{fontenc}
\usepackage[utf8]{inputenc}
\usepackage[shorthands=off,brazil]{babel}   % shorthands OFF: com babel-pt o `"` fica ativo e quebra codigo em linha
\usepackage{csquotes}

% --- Tipografia -------------------------------------------------------------
\usepackage[sfdefault]{roboto}            % Roboto como familia padrao
\usepackage[scaled=0.90]{roboto-mono}     % Roboto Mono para codigo
\usepackage{microtype}
\usepackage{textcomp}

% --- Pagina -----------------------------------------------------------------
\usepackage[
  a4paper,
  top=2.6cm, bottom=2.5cm, left=2.6cm, right=2.4cm,
  headheight=15pt, headsep=18pt, footskip=26pt
]{geometry}
\usepackage{pdflscape}
\usepackage{setspace}
\onehalfspacing

% --- Cores canônicas do Brandbook Unooria ("BLUNOORIA") -----------------------
% Regra semântica: "A cor não decora. Ela organiza."
%   - Azul Profundo (#0E1039): Estrutura, corpo, regras e títulos
%   - Lavanda (#8B85AC): Inteligência, decisões arquiteturais (ADR), IA
%   - Azul Claro (#7FA5E2): Interface, navegação, referências e notas
\usepackage{xcolor}
\definecolor{unooNavy}        {HTML}{0E1039}  % Azul Profundo — Estrutura (C100 M96 Y44 K56)
\definecolor{unooLavender}    {HTML}{8B85AC}  % Lavanda — Inteligência (C52 M48 Y15 K2)
\definecolor{unooBlue}        {HTML}{7FA5E2}  % Azul Claro — Interface (C54 M29 Y0 K0)
\definecolor{unooBlueLight}   {HTML}{B8D0F5}  % Variação clara de interface
\definecolor{unooBlueWash}    {HTML}{F0F4FC}  % Fundo suave de interface
\definecolor{unooLavenderWash}{HTML}{F3F2F8}  % Fundo suave de inteligência
\definecolor{unooLavenderSoft}{HTML}{5D577D}  % Lavanda escurecido p/ alto contraste tipográfico
\definecolor{unooInk}         {HTML}{0E1039}  % neutral-950 / texto primário ancorado na marca
\definecolor{unooInkSoft}     {HTML}{3A3C59}  % neutral-700 / texto secundário
\definecolor{unooMuted}       {HTML}{6D6E82}  % neutral-500 / legendas e metadados
\definecolor{unooGrid}        {HTML}{E3E5EE}  % Linhas de grade e bordas sutis
\definecolor{unooAxis}        {HTML}{CACDDC}  % Réguas e separadores
\definecolor{unooSurface}     {HTML}{FFFFFF}  % Fundo puro
\definecolor{unooPaper}       {HTML}{F7F8FC}  % Fundo alternado / zebra
% Tons funcionais (alertas de sistema e status):
\definecolor{unooRed}         {HTML}{C92A2A}  % danger-600
\definecolor{unooRedWash}     {HTML}{FDF2F2}
\definecolor{unooOrange}      {HTML}{D97706}  % warning-600
\definecolor{unooOrangeWash}  {HTML}{FFFBEB}
\definecolor{unooGreen}       {HTML}{2F855A}  % Status funcional aprovado
\definecolor{unooGreenWash}   {HTML}{F0FDF4}
\definecolor{unooYellow}      {HTML}{B45309}
\definecolor{unooYellowWash}  {HTML}{FEF3C7}
\definecolor{unooPink}        {HTML}{9D4EDD}

% --- Tabelas ----------------------------------------------------------------
\usepackage{booktabs}
\usepackage{tabularx}
\usepackage{longtable}
\usepackage{array}
\usepackage{multirow}
\usepackage{makecell}
\usepackage{colortbl}
\usepackage{ltablex}          % tabularx com quebra de pagina
\keepXColumns
\renewcommand{\arraystretch}{1.22}
\arrayrulecolor{unooAxis}
\setlength{\aboverulesep}{0pt}
\setlength{\belowrulesep}{0pt}
\setlength{\extrarowheight}{2pt}

% Tabelas em \small por decisao de composicao: as celulas deste documento
% misturam prosa e identificadores monoespacados longos (caminhos de arquivo,
% nomes de coluna), que NAO hifenizam. Em corpo 11 elas transbordavam a caixa;
% em \small cabem, e a tabela continua legivel porque a densidade e alta de
% proposito. `\emergencystretch` cobre o residuo.
\usepackage{etoolbox}
\AtBeginEnvironment{tabularx}{\small\setlength{\emergencystretch}{1.2em}}
\AtBeginEnvironment{longtable}{\small\setlength{\emergencystretch}{1.2em}}

\newcolumntype{L}[1]{>{\raggedright\arraybackslash}p{#1}}
\newcolumntype{C}[1]{>{\centering\arraybackslash}p{#1}}
\newcolumntype{R}[1]{>{\raggedleft\arraybackslash}p{#1}}
\newcolumntype{Y}{>{\raggedright\arraybackslash}X}

% Colunas X sem justificacao: as celulas misturam prosa e codigo monoespacado,
% que nao hifeniza — justificar produziria espacos internos muito largos.
\renewcommand{\tabularxcolumn}[1]{>{\raggedright\arraybackslash}p{#1}}

\newcommand{\thd}[1]{\textcolor{white}{\textbf{#1}}}
\newcommand{\headrow}{\rowcolor{unooNavy}}
\newcommand{\zebra}{\rowcolor{unooPaper}}

% --- Listas -----------------------------------------------------------------
\usepackage{enumitem}
\setlist[itemize]{leftmargin=1.35em, itemsep=2pt, topsep=4pt, parsep=0pt,
                  label=\textcolor{unooBlue}{\raisebox{0.18ex}{\rule{0.40em}{0.40em}}}}
\setlist[enumerate]{leftmargin=1.7em, itemsep=2pt, topsep=4pt, parsep=0pt,
                    font=\bfseries\color{unooNavy}}
\setlist[description]{leftmargin=1.4em, style=nextline,
                      font=\bfseries\color{unooNavy}}

% --- Matematica e unidades --------------------------------------------------
\usepackage{amsmath}
\usepackage{amssymb}
\usepackage{siunitx}
\sisetup{output-decimal-marker={,}, group-separator={.}, group-minimum-digits=4,
         per-mode=symbol, detect-weight=true, detect-family=true}
\DeclareSIUnit{\tokens}{tokens}
\DeclareSIUnit{\req}{req}

% --- Figuras e legendas -----------------------------------------------------
\usepackage{graphicx}
\usepackage{float}
\usepackage{adjustbox}
\usepackage[font=small, labelfont={bf,color=unooNavy}, textfont=normalfont,
            labelsep=endash, justification=raggedright, singlelinecheck=false,
            skip=8pt]{caption}
\usepackage{subcaption}
\captionsetup[table]{position=top, skip=6pt}

% Ajusta a figura a largura do texto apenas quando ela excede a caixa.
\newcommand{\fitwidth}[1]{\adjustbox{max width=\linewidth, center}{#1}}
\newcommand{\fitpage}[1]{\adjustbox{max width=\linewidth, max height=0.82\textheight, center}{#1}}

% --- Blocos de destaque (tcolorbox) ----------------------------------------
\usepackage[most]{tcolorbox}
\tcbuselibrary{skins,breakable}

\tcbset{
  unoobase/.style={
    enhanced, breakable, sharp corners,
    boxrule=0pt, leftrule=2.6pt, arc=0pt,
    left=10pt, right=10pt, top=7pt, bottom=7pt,
    fonttitle=\bfseries\small,
    before skip=10pt, after skip=10pt
  }
}

% `colbacktitle` e OBRIGATORIO aqui: sem ele o tcolorbox pinta a faixa do
% titulo com a cor de `colframe`, e o titulo (que usa a MESMA cor em
% `coltitle`) fica invisivel — texto vermelho sobre faixa vermelha. O sintoma
% e uma barra colorida solida no lugar do rotulo do bloco.
\newtcolorbox{nota}[1][]{unoobase, colback=unooBlueWash, colframe=unooBlue,
  colbacktitle=unooBlueWash, coltitle=unooNavy, titlerule=0pt, title={NOTA}, #1}
\newtcolorbox{atencao}[1][]{unoobase, colback=unooOrangeWash, colframe=unooOrange,
  colbacktitle=unooOrangeWash, coltitle=unooOrange, titlerule=0pt,
  title={ATENÇÃO}, #1}
\newtcolorbox{decisao}[1][]{unoobase, colback=unooLavenderWash, colframe=unooLavender,
  colbacktitle=unooLavenderWash, coltitle=unooNavy, titlerule=0pt,
  title={DECISÃO DE PROJETO}, #1}
\newtcolorbox{limitacao}[1][]{unoobase, colback=unooRedWash, colframe=unooRed,
  colbacktitle=unooRedWash, coltitle=unooRed, titlerule=0pt,
  title={LIMITAÇÃO CONHECIDA}, #1}
\newtcolorbox{regra}[1][]{unoobase, colback=unooRedWash, colframe=unooRed,
  colbacktitle=unooRedWash, coltitle=unooRed, titlerule=0pt,
  title={REGRA INVIOLÁVEL}, #1}
\newtcolorbox{destaque}[1][]{unoobase, colback=unooPaper, colframe=unooLavender, #1}
\newtcolorbox{especificacao}[1]{unoobase, colback=unooSurface, colframe=unooNavy,
  colbacktitle=unooSurface, coltitle=unooNavy, titlerule=0pt, title={#1}}

% --- Codigo-fonte -----------------------------------------------------------
\usepackage{listings}

\definecolor{lstBg}     {HTML}{FAFBFB}
\definecolor{lstFrame}  {HTML}{E4E9EB}
\definecolor{lstKeyword}{HTML}{2F4F7F}
\definecolor{lstString} {HTML}{0F6B52}
\definecolor{lstComment}{HTML}{617379}
\definecolor{lstEmph}   {HTML}{A9490F}

\lstdefinestyle{unoo}{
  backgroundcolor=\color{lstBg},
  basicstyle=\ttfamily\scriptsize\color{unooInk},
  keywordstyle=\color{lstKeyword}\bfseries,
  keywordstyle=[2]\color{lstEmph}\bfseries,
  keywordstyle=[3]\color{unooBlue},
  stringstyle=\color{lstString},
  commentstyle=\color{lstComment}\itshape,
  numberstyle=\ttfamily\tiny\color{unooMuted},
  identifierstyle=\color{unooInk},
  numbers=left, numbersep=9pt, stepnumber=1,
  frame=single, framesep=6pt, rulecolor=\color{lstFrame},
  xleftmargin=17pt, xrightmargin=2pt,
  breaklines=true, breakatwhitespace=false, breakindent=0pt,
  postbreak=\mbox{\textcolor{unooBlue}{$\hookrightarrow$}\space},
  showstringspaces=false, tabsize=2,
  columns=fullflexible, keepspaces=true, upquote=true,
  aboveskip=10pt, belowskip=10pt,
  literate=
    {á}{{\'a}}1 {à}{{\`a}}1 {ã}{{\~a}}1 {â}{{\^a}}1
    {é}{{\'e}}1 {ê}{{\^e}}1 {í}{{\'i}}1
    {ó}{{\'o}}1 {õ}{{\~o}}1 {ô}{{\^o}}1
    {ú}{{\'u}}1 {ü}{{\"u}}1 {ç}{{\c c}}1
    {Á}{{\'A}}1 {À}{{\`A}}1 {Ã}{{\~A}}1 {Â}{{\^A}}1
    {É}{{\'E}}1 {Ê}{{\^E}}1 {Í}{{\'I}}1
    {Ó}{{\'O}}1 {Õ}{{\~O}}1 {Ô}{{\^O}}1
    {Ú}{{\'U}}1 {Ç}{{\c C}}1
    {°}{{\textdegree}}1 {º}{{\textordmasculine}}1
    {µ}{{$\mu$}}1 {×}{{$\times$}}1 {≈}{{$\approx$}}1
    {≥}{{$\geq$}}1 {≤}{{$\leq$}}1 {≠}{{$\neq$}}1
    {→}{{$\rightarrow$}}1 {←}{{$\leftarrow$}}1 {⇒}{{$\Rightarrow$}}1
    {—}{{---}}3 {–}{{--}}2 {…}{{\ldots}}3
    {“}{{\textquotedblleft}}1 {”}{{\textquotedblright}}1
    {‘}{{\textquoteleft}}1 {’}{{\textquoteright}}1
    {ª}{{\textordfeminine}}1 {§}{{\S}}1
}
\lstset{style=unoo}

\lstdefinelanguage{tsunoo}[]{Java}{
  morekeywords={const,let,export,default,import,from,as,interface,type,await,
                async,function,return,string,number,boolean,readonly,unknown,
                never,satisfies,enum,declare,namespace,implements},
  morekeywords=[2]{z,Prisma,useEffect,postMessage},
  sensitive=true,
  morecomment=[l]{//},
  morecomment=[s]{/*}{*/},
  morestring=[b]",
  morestring=[b]',
  morestring=[b]`
}
\lstdefinelanguage{sqlunoo}[]{SQL}{
  morekeywords={POLICY,ENABLE,FORCE,ROW,LEVEL,SECURITY,USING,PARTITION,RANGE,
                CONFLICT,DO,NOTHING,JSONB,TIMESTAMPTZ,current_setting,
                set_config,ATTACH,DEFAULT,GENERATED,ALWAYS},
  sensitive=false,
  morecomment=[l]{--},
  morestring=[b]'
}
\lstdefinelanguage{yamlunoo}{
  morekeywords={true,false,null,services,image,build,command,environment,
                volumes,ports,depends_on,healthcheck,condition,runs-on,steps,
                uses,needs,jobs,name,on,with,env,if,run,concurrency,
                permissions,workflow_call,workflow_dispatch},
  sensitive=true,
  morecomment=[l]{\#},
  morestring=[b]",
  morestring=[b]'
}
\lstdefinelanguage{jsonunoo}{
  morekeywords={true,false,null},
  keywordstyle=\color{lstEmph}\bfseries,
  sensitive=true,
  morestring=[b]"
}
\lstdefinelanguage{mermaid}{
  morekeywords={flowchart,graph,sequenceDiagram,erDiagram,classDiagram,gantt,
            subgraph,end,participant,actor,alt,else,loop,par,opt,note,activate,
            deactivate,direction,TB,LR,RL,BT,Person,Person_Ext,System,System_Ext,
            SystemDb,Container,ContainerDb,ContainerQueue,Component,ComponentDb,
            Rel,Rel_D,Rel_U,Rel_L,Rel_R,System_Boundary,Container_Boundary,
            Enterprise_Boundary,Boundary,Deployment_Node,UpdateLayoutConfig,
            UpdateRelStyle,classDef,class,style,state,C4Context,C4Container,
            C4Component,C4Deployment,stateDiagram,stateDiagram-v2},
  sensitive=true,
  morecomment=[l]{\%\%},
  morestring=[b]"
}
\lstdefinelanguage{bashunoo}[]{bash}{
  morekeywords={pnpm,docker,gcloud,node,psql,tectonic,turbo,prisma,gsutil},
  sensitive=true
}

% Codigo em linha.
%
% Uso obrigatorio de \texttt e \textcolor (e nao de {\ttfamily\color{...}}):
% os comandos \text... emitem \leavevmode antes de trocar a fonte, o que faz o
% paragrafo comecar na fonte do texto corrente. Sem isso, uma celula de tabela
% que inicia com codigo teria a primeira linha de base calculada com as metricas
% da fonte monoespacada e ficaria visivelmente desalinhada das demais colunas.
\newcommand{\cd}[1]{\texttt{\textcolor{unooNavy}{#1}}}
\newcommand{\arq}[1]{\texttt{\textcolor{unooInkSoft}{#1}}}
\newcommand{\rota}[1]{\texttt{\textcolor{unooGreen}{#1}}}
\newcommand{\env}[1]{\texttt{\textcolor{lstEmph}{#1}}}
% Ponto de quebra de linha sem hifen, para caminho longo em celula estreita:
% identificador monoespacado nao hifeniza, e sem isto ele transborda a caixa.
\newcommand{\B}{\discretionary{}{}{}}

% --- Titulos ----------------------------------------------------------------
\usepackage{titlesec}
\usepackage{titletoc}

\titleformat{\chapter}[display]
  {\normalfont}
  {\raggedleft\fontsize{62}{62}\selectfont\color{unooBlueLight}\textbf{\thechapter}}
  {-10pt}
  {\raggedright\fontsize{23}{27}\selectfont\bfseries\color{unooNavy}}
  [\vspace{5pt}{\color{unooBlue}\rule{\linewidth}{2.2pt}}]
\titlespacing*{\chapter}{0pt}{-30pt}{26pt}

\titleformat{name=\chapter,numberless}[display]
  {\normalfont}
  {}
  {0pt}
  {\raggedright\fontsize{23}{27}\selectfont\bfseries\color{unooNavy}}
  [\vspace{5pt}{\color{unooBlue}\rule{\linewidth}{2.2pt}}]

\titleformat{\section}
  {\normalfont\Large\bfseries\color{unooNavy}}{\thesection}{0.7em}{}
\titleformat{\subsection}
  {\normalfont\large\bfseries\color{unooInk}}{\thesubsection}{0.7em}{}
\titleformat{\subsubsection}
  {\normalfont\normalsize\bfseries\color{unooInkSoft}}{\thesubsubsection}{0.7em}{}
\titleformat{\paragraph}[runin]
  {\normalfont\normalsize\bfseries\color{unooNavy}}{}{0pt}{}[\;---\;]

\titlespacing*{\section}{0pt}{20pt}{8pt}
\titlespacing*{\subsection}{0pt}{15pt}{6pt}
\titlespacing*{\subsubsection}{0pt}{12pt}{5pt}
\titlespacing*{\paragraph}{0pt}{9pt}{0.4em}

\setcounter{secnumdepth}{3}
\setcounter{tocdepth}{2}

% --- Sumario ----------------------------------------------------------------
\titlecontents{chapter}[0em]
  {\addvspace{9pt}\bfseries\color{unooNavy}}
  {\contentslabel[\thecontentslabel]{1.5em}}
  {}
  {\hfill\contentspage}
\titlecontents{section}[1.7em]
  {\addvspace{1.5pt}}
  {\contentslabel{2.2em}}
  {}
  {\titlerule*[0.7pc]{\color{unooAxis}.}\contentspage}
\titlecontents{subsection}[4.0em]
  {\addvspace{0.5pt}\small\color{unooInkSoft}}
  {\contentslabel{2.8em}}
  {}
  {\titlerule*[0.7pc]{\color{unooGrid}.}\contentspage}

% --- Cabecalho e rodape -----------------------------------------------------
\usepackage{fancyhdr}
\usepackage{lastpage}

\newcommand{\docid}{PU-DOC-001}
\newcommand{\docversao}{1.2}

\renewcommand{\headrule}{{\color{unooGrid}\hrule height \headrulewidth}}
\renewcommand{\footrule}{{\color{unooGrid}\hrule height \footrulewidth}}

\fancypagestyle{unoomain}{%
  \fancyhf{}
  \fancyhead[L]{\footnotesize\color{unooInkSoft}\nouppercase{\leftmark}}
  \fancyhead[R]{\raisebox{-0.15ex}{\includegraphics[height=1.8ex]{figuras/marca/simbolo_azul_profundo.png}}\hspace{4pt}\footnotesize\bfseries\color{unooNavy}Unooria}
  \fancyfoot[L]{\footnotesize\color{unooMuted}\docid{} \textperiodcentered{} v\docversao}
  \fancyfoot[R]{\footnotesize\color{unooInkSoft}\thepage\,/\,\pageref{LastPage}}
  \renewcommand{\headrulewidth}{0.5pt}
  \renewcommand{\footrulewidth}{0.5pt}
}
% Pre-textual: numeracao romana, sem total de paginas (o total do
% \pageref{LastPage} refere-se a numeracao arabe do corpo do documento).
\fancypagestyle{unoofront}{%
  \fancyhf{}
  \fancyhead[L]{\footnotesize\color{unooInkSoft}\nouppercase{\leftmark}}
  \fancyhead[R]{\raisebox{-0.15ex}{\includegraphics[height=1.8ex]{figuras/marca/simbolo_azul_profundo.png}}\hspace{4pt}\footnotesize\bfseries\color{unooNavy}Unooria}
  \fancyfoot[L]{\footnotesize\color{unooMuted}\docid{} \textperiodcentered{} v\docversao}
  \fancyfoot[R]{\footnotesize\color{unooInkSoft}\thepage}
  \renewcommand{\headrulewidth}{0.5pt}
  \renewcommand{\footrulewidth}{0.5pt}
}
\fancypagestyle{plain}{%
  \fancyhf{}
  \fancyfoot[R]{\footnotesize\color{unooInkSoft}\thepage}
  \renewcommand{\headrulewidth}{0pt}
  \renewcommand{\footrulewidth}{0pt}
}
\pagestyle{unoofront}
\renewcommand{\chaptermark}[1]{\markboth{\thechapter\quad #1}{}}

% --- Referencias cruzadas e PDF --------------------------------------------
\usepackage[
  colorlinks=true,
  linkcolor=unooNavy,
  citecolor=unooBlue,
  urlcolor=unooBlue,
  filecolor=unooBlue,
  bookmarksnumbered=true,
  bookmarksopen=true,
  bookmarksopenlevel=1,
  pdfdisplaydoctitle=true
]{hyperref}
\usepackage{bookmark}
\usepackage{url}
\urlstyle{sf}

% --- Bibliografia (ABNT NBR 6023) ------------------------------------------
%
% `abntex2cite` + BibTeX, e nao biblatex+biber: o pipeline de compilacao
% (Tectonic) chama o BibTeX classico automaticamente, sem ferramenta externa.
% A saida sai formatada em ABNT do mesmo jeito.
\usepackage[alf]{abntex2cite}

% --- Semantica de referencia do projeto ------------------------------------
\newcommand{\rf}[1]{\textbf{\textcolor{unooNavy}{#1}}}
\newcommand{\adrr}[1]{\textbf{\textcolor{unooLavenderSoft}{#1}}}
\newcommand{\ucr}[1]{\textbf{\textcolor{unooBlue}{#1}}}
\newcommand{\drr}[1]{\textbf{\textcolor{unooOrange}{#1}}}
\newcommand{\yes}{\textcolor{unooGreen}{\textbf{\checkmark}}}
\newcommand{\no}{\textcolor{unooOrange}{\textbf{--}}}

% Setas: a Roboto NAO tem U+2192 nem U+21D2 (medido — o XeTeX avisa
% "Missing character" e o glifo simplesmente nao e desenhado). Em prosa,
% sempre estes comandos; nunca a seta literal.
\newcommand{\seta}{\ensuremath{\rightarrow}}
\newcommand{\implica}{\ensuremath{\Rightarrow}}
\newcommand{\setabi}{\ensuremath{\leftrightarrow}}

% Espacamento de paragrafo
\setlength{\parindent}{0pt}
\setlength{\parskip}{7pt}

\widowpenalty=10000
\clubpenalty=10000
\flushbottom

% ============================================================================
%  diagramas.tex — biblioteca TikZ do documento
%
%  Todos os diagramas do documento (C4, UML de casos de uso, atividades,
%  sequencia, entidade-relacionamento, classes e infraestrutura) sao vetoriais
%  e desenhados nativamente em TikZ. Isso garante tipografia unica (Roboto),
%  nitidez em qualquer zoom e ausencia de dependencia externa de renderizacao.
%  O codigo Mermaid equivalente de cada diagrama esta no Apendice A.
% ============================================================================

\usepackage{tikz}
\usetikzlibrary{
  arrows.meta, positioning, fit, backgrounds, calc, shapes.geometric,
  shapes.symbols, shapes.multipart, shapes.misc, decorations.pathreplacing,
  decorations.pathmorphing, decorations.markings, patterns, patterns.meta,
  matrix, chains, scopes, quotes, shadows.blur, math, intersections
}

\usepackage{pgfplots}
\pgfplotsset{compat=1.18}
\usepgfplotslibrary{groupplots}

% Numeracao dos graficos na convencao brasileira (virgula decimal), coerente
% com o restante do documento.
\pgfkeys{/pgf/number format/.cd, use comma, 1000 sep={.}}

% Estilo comum dos graficos quantitativos: eixos recessivos, grade sutil,
% tipografia do documento. Mesma paleta validada da interface.
\pgfplotsset{
  unooplot/.style={
    width=\linewidth, height=6cm,
    axis line style={draw=unooAxis, line width=0.6pt},
    tick label style={font=\scriptsize\color{unooInkSoft}},
    label style={font=\small\color{unooInkSoft}},
    title style={font=\small\bfseries\color{unooNavy}, align=left},
    legend style={font=\scriptsize, draw=unooGrid, fill=unooSurface,
                  /tikz/every even column/.append style={column sep=6pt}},
    grid=major, grid style={draw=unooGrid, line width=0.5pt},
    major tick style={draw=unooAxis},
    axis background/.style={fill=unooSurface},
    every axis plot post/.append style={line width=0.9pt},
    enlarge x limits=0.06,
  }
}

% ---------------------------------------------------------------------------
%  Base comum
% ---------------------------------------------------------------------------
\tikzset{
  >={Stealth[length=5.4pt, width=4.2pt]},
  every node/.append style={align=center, inner sep=4pt},
  cxlabel/.style={font=\scriptsize\color{unooInkSoft}, fill=unooSurface,
                  inner sep=1.6pt, rounded corners=1pt},
  cxtitle/.style={font=\small\bfseries\color{unooNavy}},
  legenda/.style={font=\scriptsize\color{unooMuted}},
}

% ---------------------------------------------------------------------------
%  C4 model — notacao de Simon Brown adaptada a paleta do projeto
% ---------------------------------------------------------------------------
\tikzset{
  % --- elementos -----------------------------------------------------------
  c4box/.style={
    rectangle, rounded corners=2.5pt, draw, line width=0.8pt,
    text width=32mm, minimum height=15mm, font=\footnotesize,
    inner sep=5pt
  },
  c4person/.style={
    c4box, fill=unooNavy, draw=unooNavy!85!black, text=white,
    rounded corners=7pt, text width=28mm, minimum height=16mm
  },
  c4system/.style={
    c4box, fill=unooBlue, draw=unooNavy, text=white
  },
  c4systemext/.style={
    c4box, fill=unooMuted, draw=unooInkSoft, text=white
  },
  c4container/.style={
    c4box, fill=unooBlue!88, draw=unooNavy, text=white
  },
  c4containerext/.style={
    c4box, fill=unooMuted!85, draw=unooInkSoft, text=white
  },
  c4component/.style={
    c4box, fill=unooBlueLight!45, draw=unooBlue, text=unooInk, text width=30mm,
    minimum height=13mm
  },
  c4db/.style={
    cylinder, shape border rotate=90, aspect=0.28, draw, line width=0.8pt,
    fill=unooNavy, draw=unooNavy!80!black, text=white,
    text width=25mm, minimum height=17mm, font=\footnotesize, inner sep=3pt
  },
  c4queue/.style={
    c4box, fill=unooYellow!85, draw=unooYellow!60!black, text=unooInk,
    rounded corners=8pt
  },
  c4artifact/.style={
    c4box, fill=unooGreen!80, draw=unooGreen!60!black, text=white
  },
  % --- fronteiras ----------------------------------------------------------
  c4boundary/.style={
    draw=unooAxis, dashed, dash pattern=on 3pt off 2.4pt, line width=0.9pt,
    rounded corners=4pt, inner sep=9pt
  },
  c4boundarylbl/.style={
    font=\scriptsize\bfseries\color{unooInkSoft}, anchor=north west,
    inner sep=2pt
  },
  c4zone/.style={
    c4boundary, draw=unooBlue!55, fill=unooBlue!4
  },
  c4zoneext/.style={
    c4boundary, draw=unooMuted, fill=unooPaper
  },
  % --- relacoes ------------------------------------------------------------
  c4rel/.style={-{Stealth[length=5pt,width=4pt]}, line width=0.85pt,
                draw=unooInkSoft},
  c4relbi/.style={c4rel, {Stealth[length=5pt,width=4pt]}-{Stealth[length=5pt,width=4pt]}},
  c4reldash/.style={c4rel, dashed, dash pattern=on 3pt off 2pt},
  c4relhi/.style={c4rel, draw=unooOrange, line width=1.1pt},
}

% Rotulo tecnologico dentro de um elemento C4
\newcommand{\tech}[1]{\\[1.5pt]{\scriptsize\itshape [#1]}}
% Descricao curta dentro de um elemento C4
\newcommand{\desc}[1]{\\[2pt]{\scriptsize #1}}

% ---------------------------------------------------------------------------
%  UML — casos de uso
% ---------------------------------------------------------------------------
\tikzset{
  ucase/.style={
    ellipse, draw=unooBlue, line width=0.85pt, fill=unooBlueWash,
    text=unooInk, font=\scriptsize, minimum width=34mm, minimum height=11mm,
    inner sep=2.5pt, text width=31mm
  },
  ucasesec/.style={ucase, fill=unooSurface, draw=unooAxis},
  ucsystem/.style={
    draw=unooNavy, line width=1pt, rounded corners=3pt, fill=unooBlue!3,
    inner sep=11pt
  },
  ucrel/.style={line width=0.8pt, draw=unooInkSoft},
  ucinclude/.style={
    -{Stealth[length=5pt,width=4pt]}, dashed, dash pattern=on 3pt off 2pt,
    line width=0.75pt, draw=unooGreen
  },
  ucextend/.style={
    -{Stealth[length=5pt,width=4pt]}, dashed, dash pattern=on 3pt off 2pt,
    line width=0.75pt, draw=unooOrange
  },
  ucgen/.style={-{Triangle[open, length=6pt, width=6pt]}, line width=0.8pt,
                draw=unooInkSoft},
  ucnote/.style={font=\scriptsize\color{unooInkSoft}},
}

% Ator UML (boneco de palito) — #1 nome do no, #2 posicao, #3 rotulo
\newcommand{\ator}[3]{%
  \begin{scope}[shift={#2}]
    \draw[line width=0.9pt, draw=unooNavy] (0,0.30) circle (0.155);
    \draw[line width=0.9pt, draw=unooNavy] (0,0.145) -- (0,-0.24);
    \draw[line width=0.9pt, draw=unooNavy] (-0.245,-0.01) -- (0.245,-0.01);
    \draw[line width=0.9pt, draw=unooNavy] (0,-0.24) -- (-0.20,-0.60);
    \draw[line width=0.9pt, draw=unooNavy] (0,-0.24) -- (0.20,-0.60);
    \node[font=\scriptsize\bfseries\color{unooNavy}, text width=24mm,
          anchor=north] at (0,-0.68) {#3};
    \node[inner sep=0pt, minimum width=13mm, minimum height=19mm] (#1) at (0,-0.15) {};
  \end{scope}
}

% Ator secundario (sistema externo) — retangulo com estereotipo
\newcommand{\atorsys}[3]{%
  \node[draw=unooInkSoft, line width=0.85pt, fill=unooPaper, rounded corners=2pt,
        font=\scriptsize, text width=22mm, minimum height=12mm] (#1) at #2
        {{\scriptsize$\ll$ator$\gg$}\\\textbf{#3}};
}

% ---------------------------------------------------------------------------
%  UML — diagramas de atividades
% ---------------------------------------------------------------------------
\tikzset{
  actstart/.style={
    circle, fill=unooInk, draw=unooInk, minimum size=4.6mm, inner sep=0pt
  },
  actend/.style={
    circle, fill=unooInk, draw=unooInk, minimum size=4.6mm, inner sep=0pt,
    double, double distance=1.1pt, outer sep=0pt
  },
  action/.style={
    rectangle, rounded corners=4pt, draw=unooBlue, line width=0.85pt,
    fill=unooBlueWash, text=unooInk, font=\scriptsize, text width=36mm,
    minimum height=10mm, inner sep=4pt
  },
  actionsys/.style={action, fill=unooSurface, draw=unooAxis},
  actionerr/.style={action, fill=unooOrangeWash, draw=unooOrange},
  actionok/.style={action, fill=unooGreenWash, draw=unooGreen},
  decision/.style={
    diamond, draw=unooNavy, line width=0.85pt, fill=unooYellowWash,
    aspect=1.9, font=\scriptsize, text width=26mm, inner sep=1pt
  },
  bar/.style={
    rectangle, fill=unooInk, minimum width=34mm, minimum height=1.5mm,
    inner sep=0pt
  },
  barv/.style={
    rectangle, fill=unooInk, minimum width=1.5mm, minimum height=30mm,
    inner sep=0pt
  },
  flow/.style={-{Stealth[length=5pt,width=4pt]}, line width=0.85pt,
               draw=unooInkSoft},
  flowlbl/.style={font=\scriptsize\color{unooInkSoft}, fill=unooSurface,
                  inner sep=1.5pt},
  lane/.style={
    draw=unooAxis, line width=0.7pt, fill=unooSurface
  },
  lanehdr/.style={
    draw=unooAxis, line width=0.7pt, fill=unooNavy, text=white,
    font=\scriptsize\bfseries
  },
  actnote/.style={
    draw=unooYellow!70!black, fill=unooYellowWash, line width=0.6pt,
    font=\scriptsize\color{unooInkSoft}, text width=34mm, align=left,
    rounded corners=0pt, inner sep=4pt
  },
  actnoteline/.style={dashed, dash pattern=on 2pt off 1.6pt, draw=unooYellow!70!black},
  objflow/.style={
    rectangle, draw=unooGreen, line width=0.7pt, fill=unooGreenWash,
    font=\scriptsize, text width=30mm, inner sep=3.5pt
  },
}

% ---------------------------------------------------------------------------
%  UML — diagramas de sequencia
% ---------------------------------------------------------------------------
\tikzset{
  seqhead/.style={
    rectangle, rounded corners=2pt, draw=unooNavy, line width=0.8pt,
    fill=unooNavy, text=white, font=\scriptsize, text width=26mm,
    minimum height=9mm, inner sep=3pt
  },
  seqheadext/.style={seqhead, fill=unooMuted, draw=unooInkSoft},
  seqheadact/.style={seqhead, fill=unooBlue, draw=unooNavy},
  seqline/.style={dashed, dash pattern=on 2.6pt off 2.2pt, draw=unooAxis,
                  line width=0.7pt},
  seqmsg/.style={-{Stealth[length=5pt,width=4pt]}, line width=0.85pt,
                 draw=unooInkSoft},
  seqret/.style={seqmsg, dashed, dash pattern=on 2.6pt off 2pt,
                 draw=unooInkSoft},
  seqself/.style={seqmsg, rounded corners=2pt},
  seqlbl/.style={font=\scriptsize\color{unooInk}, fill=unooSurface,
                 inner sep=1.6pt},
  seqact/.style={fill=unooBlueLight!60, draw=unooBlue, line width=0.5pt,
                 minimum width=2.2mm, inner sep=0pt},
  seqframe/.style={draw=unooAxis, line width=0.7pt, rounded corners=0pt},
  seqframelbl/.style={fill=unooNavy, text=white, font=\scriptsize\bfseries,
                      inner sep=2.2pt, anchor=north west},
  seqcond/.style={font=\scriptsize\itshape\color{unooInkSoft}, anchor=west},
}

% Raia de sequencia: #1 = x, #2 = profundidade (y final, negativo),
% #3 = titulo, #4 = tecnologia
\newcommand{\raia}[4]{%
  \node[seqhead] at (#1,0) {\textbf{#3}\\[0.5pt]{\scriptsize\itshape #4}};
  \draw[seqline] (#1,-0.52) -- (#1,#2);
}
\newcommand{\raiaext}[4]{%
  \node[seqheadext] at (#1,0) {\textbf{#3}\\[0.5pt]{\scriptsize\itshape #4}};
  \draw[seqline] (#1,-0.52) -- (#1,#2);
}
% Mensagem: #1 = x origem, #2 = x destino, #3 = y, #4 = rotulo
\newcommand{\msgs}[4]{%
  \draw[seqmsg] (#1,#3) -- (#2,#3) node[seqlbl, midway, above=0.4pt]{#4};
}
% Retorno (tracejado)
\newcommand{\msgr}[4]{%
  \draw[seqret] (#1,#3) -- (#2,#3) node[seqlbl, midway, above=0.4pt]{#4};
}
% Auto-mensagem: #1 = x, #2 = y, #3 = rotulo
\newcommand{\msgself}[3]{%
  \draw[seqself] (#1,#2) -- ++(0.62,0) -- ++(0,-0.42) -- ++(-0.62,0);
  \node[seqlbl, anchor=west] at (#1+0.72,#2-0.21) {#3};
}

% ---------------------------------------------------------------------------
%  Entidade-relacionamento e classes
% ---------------------------------------------------------------------------
\tikzset{
  entidade/.style={
    rectangle split, rectangle split parts=2, draw=unooNavy,
    line width=0.85pt, rectangle split part fill={unooNavy, unooSurface},
    font=\scriptsize, text width=37mm, inner sep=4pt, rounded corners=1pt,
    every text node part/.style={align=center},
    every two node part/.style={align=left}
  },
  classe/.style={
    rectangle split, rectangle split parts=3, draw=unooNavy,
    line width=0.85pt, rectangle split part fill={unooNavy, unooSurface, unooPaper},
    font=\scriptsize, text width=44mm, inner sep=4pt, rounded corners=1pt,
    every text node part/.style={align=center},
    every two node part/.style={align=left},
    every three node part/.style={align=left}
  },
  erel/.style={line width=0.85pt, draw=unooInkSoft},
  ercard/.style={font=\scriptsize\color{unooNavy}, fill=unooSurface,
                 inner sep=1.4pt},
  hered/.style={-{Triangle[open, length=6.5pt, width=6.5pt]}, line width=0.85pt,
                draw=unooInkSoft},
  compos/.style={-{Diamond[open, length=6pt, width=4.5pt]}, line width=0.85pt,
                 draw=unooInkSoft},
  depend/.style={-{Stealth[length=5pt,width=4pt]}, dashed,
                 dash pattern=on 3pt off 2pt, line width=0.75pt,
                 draw=unooInkSoft},
}

% Titulo branco em negrito para a primeira particao de entidade/classe
\newcommand{\etitulo}[1]{\textcolor{white}{\textbf{#1}}}
\newcommand{\pk}[1]{\textcolor{unooNavy}{\textbf{#1}}}
\newcommand{\fk}[1]{\textcolor{unooOrange}{#1}}

% ---------------------------------------------------------------------------
%  Infraestrutura / implantacao
% ---------------------------------------------------------------------------
\tikzset{
  nuvem/.style={
    draw=unooBlue, line width=1pt, dash pattern=on 4pt off 2.6pt,
    rounded corners=6pt, fill=unooBlue!3, inner sep=11pt
  },
  nuvemlbl/.style={
    font=\scriptsize\bfseries\color{unooNavy}, anchor=north west,
    fill=unooSurface, inner sep=2.4pt
  },
  regiao/.style={
    draw=unooAxis, line width=0.8pt, dash pattern=on 2.6pt off 2pt,
    rounded corners=4pt, fill=unooSurface, inner sep=8pt
  },
  servico/.style={
    rectangle, rounded corners=2.5pt, draw=unooNavy, line width=0.85pt,
    fill=unooBlue!88, text=white, font=\scriptsize, text width=30mm,
    minimum height=13mm, inner sep=4pt
  },
  servicoger/.style={servico, fill=unooNavy},
  servicobroker/.style={servico, fill=unooYellow!85, text=unooInk,
                        draw=unooYellow!55!black},
  volume/.style={
    cylinder, shape border rotate=90, aspect=0.30, draw=unooGreen!60!black,
    line width=0.85pt, fill=unooGreen!80, text=white, font=\scriptsize,
    text width=22mm, minimum height=15mm, inner sep=3pt
  },
  externo/.style={
    rectangle, rounded corners=2.5pt, draw=unooInkSoft, line width=0.85pt,
    fill=unooMuted, text=white, font=\scriptsize, text width=30mm,
    minimum height=13mm, inner sep=4pt
  },
  internet/.style={
    cloud, cloud puffs=13, cloud puff arc=110, draw=unooAxis, line width=0.85pt,
    fill=unooPaper, text=unooInkSoft, font=\scriptsize, minimum width=26mm,
    minimum height=15mm, aspect=1.6
  },
  netpriv/.style={-{Stealth[length=5pt,width=4pt]}, line width=0.9pt,
                  draw=unooGreen!70!black},
  netpub/.style={-{Stealth[length=5pt,width=4pt]}, line width=0.9pt,
                 draw=unooOrange},
  netbi/.style={{Stealth[length=5pt,width=4pt]}-{Stealth[length=5pt,width=4pt]},
                line width=0.9pt, draw=unooInkSoft},
  portlbl/.style={font=\scriptsize\ttfamily\color{unooInkSoft},
                  fill=unooSurface, inner sep=1.5pt},
}

% ---------------------------------------------------------------------------
%  Pipeline / fluxo de dados
% ---------------------------------------------------------------------------
\tikzset{
  etapa/.style={
    rectangle, rounded corners=2.5pt, draw=unooBlue, line width=0.85pt,
    fill=unooBlueWash, font=\scriptsize, text width=27mm, minimum height=13mm,
    inner sep=4pt
  },
  etapaio/.style={
    etapa, fill=unooPaper, draw=unooAxis
  },
  etapaout/.style={etapa, fill=unooGreenWash, draw=unooGreen},
  etapaerr/.style={etapa, fill=unooOrangeWash, draw=unooOrange},
  fita/.style={
    signal, signal to=east, signal from=west, draw=unooBlue,
    line width=0.85pt, fill=unooBlueWash, font=\scriptsize, text width=24mm,
    minimum height=12mm
  },
  pipe/.style={-{Stealth[length=5.4pt,width=4.4pt]}, line width=1pt,
               draw=unooBlue},
  pipelbl/.style={font=\scriptsize\color{unooInkSoft}, fill=unooSurface,
                  inner sep=1.5pt},
}

% ---------------------------------------------------------------------------
%  Legenda padronizada de diagrama
% ---------------------------------------------------------------------------
\newcommand{\amostra}[1]{%
  \tikz[baseline=-0.55ex]{\node[#1, text width=0pt, minimum width=3.6mm,
    minimum height=3.0mm, inner sep=0pt, font=\tiny]{};}%
}
\newenvironment{legendadiag}
  {\par\vspace{2pt}\begingroup\footnotesize\color{unooInkSoft}%
   \setlength{\parskip}{2pt}\textbf{Legenda:}\space}
  {\endgroup\par\vspace{2pt}}

% ---------------------------------------------------------------------------
%  Acrescimos deste documento (nao vieram da biblioteca de origem)
% ---------------------------------------------------------------------------
\tikzset{
  % Bundle estatico servido pela CDN — nao e um servico, e um artefato.
  bundle/.style={
    rectangle, draw=unooGreen!60!black, line width=0.85pt, fill=unooGreen!80,
    text=white, font=\scriptsize, text width=28mm, minimum height=13mm,
    inner sep=4pt, rounded corners=1pt
  },
  % Fronteira do iframe: a linha que separa o chrome do hub da ficcao da
  % experiencia. Tracejada e grossa de proposito — e fronteira de seguranca.
  iframe/.style={
    draw=unooOrange, line width=1.1pt, dash pattern=on 4pt off 2.4pt,
    rounded corners=3pt, fill=unooOrangeWash!45, inner sep=9pt
  },
  iframelbl/.style={
    font=\scriptsize\bfseries\color{unooOrange}, anchor=north west,
    fill=unooSurface, inner sep=2.2pt
  },
  % Maquina de estados
  estado/.style={
    rectangle, rounded corners=6pt, draw=unooBlue, line width=0.85pt,
    fill=unooBlueWash, font=\scriptsize, text width=24mm, minimum height=9mm,
    inner sep=3pt
  },
  estadoini/.style={estado, fill=unooNavy, text=white, draw=unooNavy},
  estadofim/.style={estado, fill=unooGreenWash, draw=unooGreen},
  estadoerr/.style={estado, fill=unooRedWash, draw=unooRed},
  trans/.style={-{Stealth[length=5pt,width=4pt]}, line width=0.8pt,
                draw=unooInkSoft},
  translbl/.style={font=\scriptsize\color{unooInkSoft}, fill=unooSurface,
                   inner sep=1.4pt},
  % Camada horizontal (diagramas de camadas)
  camada/.style={
    rectangle, draw=unooAxis, line width=0.7pt, fill=unooSurface,
    minimum height=13mm, inner sep=5pt, font=\scriptsize
  },
  camadalbl/.style={font=\scriptsize\bfseries\color{unooNavy}},
  % Tabela do banco, com prefixo por experiencia
  tabela/.style={
    rectangle split, rectangle split parts=2, draw=unooNavy,
    line width=0.8pt, rectangle split part fill={unooNavy, unooSurface},
    font=\scriptsize, text width=34mm, inner sep=3.5pt, rounded corners=1pt,
    every text node part/.style={align=center},
    every two node part/.style={align=left}
  },
  tabelaplat/.style={tabela, rectangle split part fill={unooBlue, unooSurface}},
  tabelacat/.style={tabela, rectangle split part fill={unooMuted, unooPaper}},
}

% Rotulo de prefixo de tabela por experiencia
\newcommand{\pref}[1]{{\scriptsize\ttfamily #1}}

% ---------------------------------------------------------------------------
%  UML 2.5.1 --- os tipos que a biblioteca de origem nao tinha (Secao
%  "Modelagem UML", secoes/05b-modelagem-uml.tex; ADR-122)
%
%  Pacote, componente com interfaces fornecidas e requeridas, no de
%  implantacao, artefato, objeto, estrutura composta (parte e porta), perfil,
%  fragmento de interacao, comunicacao e tempo. A notacao segue a
%  especificacao da OMG (UML 2.5.1, formal/17-12-05); a cor segue os tokens
%  unoo* como o resto do documento.
% ---------------------------------------------------------------------------
\newcommand{\estereo}[1]{{\guillemotleft}#1{\guillemotright}}

\tikzset{
  % --- relacoes de classe que faltavam ------------------------------------
  associa/.style={line width=0.85pt, draw=unooInkSoft},
  navega/.style={associa, -{Stealth[length=5pt,width=4pt]}},
  agrega/.style={associa, {Diamond[open, length=7.5pt, width=5pt]}-},
  compoe/.style={associa, {Diamond[length=7.5pt, width=5pt, fill=unooInkSoft]}-},
  realiza/.style={-{Triangle[open, length=6.5pt, width=6.5pt]}, dashed,
                  dash pattern=on 3pt off 2pt, line width=0.8pt, draw=unooInkSoft},
  % extensao de perfil: seta cheia, que nao se confunde com a heranca (aberta)
  estende/.style={-{Triangle[length=6.5pt, width=6.5pt, fill=unooInkSoft]},
                  line width=0.85pt, draw=unooInkSoft},
  multip/.style={font=\scriptsize\color{unooNavy}, inner sep=1.2pt},
  papelassoc/.style={font=\scriptsize\itshape\color{unooInkSoft}, inner sep=1.2pt},
  % --- classificadores ----------------------------------------------------
  interface/.style={classe, draw=unooLavenderSoft,
    rectangle split part fill={unooLavenderSoft, unooSurface, unooPaper}},
  enumeracao/.style={entidade, draw=unooInkSoft, text width=30mm,
    rectangle split part fill={unooMuted, unooSurface}},
  metaclasse/.style={rectangle, draw=unooInkSoft, line width=0.8pt,
    fill=unooPaper, font=\scriptsize, text width=28mm, minimum height=11mm,
    inner sep=3pt},
  % --- objeto: nome sublinhado, instancia e nao tipo -----------------------
  objeto/.style={rectangle split, rectangle split parts=2, draw=unooBlue,
    line width=0.85pt, rectangle split part fill={unooBlueLight!55, unooSurface},
    font=\scriptsize, text width=40mm, inner sep=3.5pt,
    every text node part/.style={align=center},
    every two node part/.style={align=left}},
  objetoum/.style={rectangle, draw=unooBlue, line width=0.85pt,
    fill=unooBlueLight!45, font=\scriptsize, text width=30mm,
    minimum height=10mm, inner sep=3pt},
  objetolink/.style={line width=0.8pt, draw=unooInkSoft},
  % --- pacote --------------------------------------------------------------
  pacotecorpo/.style={draw=unooNavy, line width=0.85pt, fill=unooBlueWash,
    inner sep=0pt},
  pacoteaba/.style={draw=unooNavy, line width=0.85pt, fill=unooNavy,
    text=white, font=\scriptsize\bfseries, inner sep=2.4pt,
    minimum height=4.6mm},
  pacotedep/.style={-{Stealth[length=5pt,width=4pt]}, dashed,
    dash pattern=on 3pt off 2pt, line width=0.8pt, draw=unooInkSoft},
  % --- componente: retangulo com o icone de componente no canto -------------
  componente/.style={rectangle, draw=unooNavy, line width=0.85pt,
    fill=unooBlueWash, font=\scriptsize, text width=30mm, minimum height=15mm,
    inner sep=4pt,
    path picture={
      \coordinate (icc) at ([xshift=-3.0mm, yshift=-3.1mm]path picture bounding box.north east);
      \draw[draw=unooNavy, line width=0.45pt, fill=unooSurface]
        ([xshift=-1.3mm, yshift=-1.8mm]icc) rectangle ([xshift=1.6mm, yshift=1.8mm]icc);
      \draw[draw=unooNavy, line width=0.45pt, fill=unooSurface]
        ([xshift=-2.1mm, yshift=0.5mm]icc) rectangle ([xshift=-0.5mm, yshift=1.2mm]icc);
      \draw[draw=unooNavy, line width=0.45pt, fill=unooSurface]
        ([xshift=-2.1mm, yshift=-1.2mm]icc) rectangle ([xshift=-0.5mm, yshift=-0.5mm]icc);
    }},
  componenteext/.style={componente, fill=unooPaper, draw=unooInkSoft},
  ifacelbl/.style={font=\scriptsize\color{unooInkSoft}, inner sep=1.4pt},
  % --- implantacao ---------------------------------------------------------
  nodocorpo/.style={draw=unooNavy, line width=0.85pt, fill=unooSurface,
    inner sep=0pt},
  nodolbl/.style={font=\scriptsize\color{unooNavy}, align=left, inner sep=3pt},
  artefato/.style={rectangle, draw=unooGreen!60!black, line width=0.8pt,
    fill=unooGreenWash, font=\scriptsize, text width=30mm, minimum height=11mm,
    inner sep=3.5pt,
    path picture={
      \coordinate (ica) at ([xshift=-2.9mm, yshift=-2.9mm]path picture bounding box.north east);
      \draw[draw=unooGreen!60!black, line width=0.45pt, fill=unooSurface]
        ([xshift=-1.3mm, yshift=1.7mm]ica) -- ([xshift=-1.3mm, yshift=-1.7mm]ica)
        -- ([xshift=1.3mm, yshift=-1.7mm]ica) -- ([xshift=1.3mm, yshift=0.8mm]ica)
        -- ([xshift=0.4mm, yshift=1.7mm]ica) -- cycle;
      \draw[draw=unooGreen!60!black, line width=0.45pt]
        ([xshift=0.4mm, yshift=1.7mm]ica) -- ([xshift=0.4mm, yshift=0.8mm]ica)
        -- ([xshift=1.3mm, yshift=0.8mm]ica);
    }},
  implanta/.style={-{Stealth[length=5pt,width=4pt]}, dashed,
    dash pattern=on 3pt off 2pt, line width=0.75pt, draw=unooGreen!60!black},
  caminhocom/.style={line width=0.9pt, draw=unooInkSoft},
  % --- estrutura composta: parte, porta e conector -------------------------
  parte/.style={rectangle, draw=unooNavy, line width=0.85pt, fill=unooSurface,
    font=\scriptsize, text width=28mm, minimum height=11mm, inner sep=3pt},
  partemult/.style={parte, fill=unooBlueWash},
  porta/.style={rectangle, draw=unooNavy, line width=0.8pt, fill=unooSurface,
    minimum size=2.6mm, inner sep=0pt},
  portaborda/.style={porta, fill=unooOrange!85, draw=unooOrange!70!black},
  conector/.style={line width=0.85pt, draw=unooInkSoft},
  % --- interacao: operador de fragmento e uso de interacao (ref) -----------
  frameop/.style={chamfered rectangle, chamfered rectangle corners=south east,
    chamfered rectangle xsep=1.4mm, chamfered rectangle ysep=1.2mm,
    fill=unooNavy, draw=unooNavy, text=white, font=\scriptsize\bfseries,
    inner sep=1.8pt, anchor=north west},
  guarda/.style={font=\scriptsize\itshape\color{unooInkSoft}, anchor=west,
    inner sep=1.4pt},
  refint/.style={rectangle, draw=unooNavy, line width=0.85pt, fill=unooSurface,
    font=\scriptsize, text width=40mm, minimum height=11mm, inner sep=3pt},
  % --- comunicacao ---------------------------------------------------------
  msgcomlbl/.style={font=\scriptsize\color{unooInk}, fill=unooSurface,
    inner sep=1.4pt},
  % --- maquina de estados: pseudoestados e estado composto -----------------
  escolha/.style={diamond, draw=unooNavy, line width=0.85pt, fill=unooSurface,
    minimum size=6mm, inner sep=0pt},
  estadocomposto/.style={draw=unooBlue, line width=0.9pt, rounded corners=8pt,
    fill=unooBlue!4},
  % --- tempo ---------------------------------------------------------------
  tempofaixa/.style={draw=unooAxis, line width=0.7pt, fill=unooSurface},
  tempolinha/.style={line width=1.2pt, draw=unooBlue!70!unooNavy},
  tempoestado/.style={font=\scriptsize\color{unooInkSoft}, anchor=east,
    inner sep=1.5pt},
  tempotick/.style={font=\scriptsize\color{unooMuted}, inner sep=1.5pt},
  tempomarca/.style={dashed, dash pattern=on 2pt off 1.6pt, draw=unooAxis,
    line width=0.6pt},
  restricao/.style={font=\scriptsize\color{unooNavy}, fill=unooSurface,
    inner sep=1.4pt},
}

% Titulo de classificador com estereotipo: #1 estereotipo, #2 nome (branco,
% para a primeira particao escura de classe/entidade/interface)
\newcommand{\ctitulo}[2]{\textcolor{white}{{\scriptsize\estereo{#1}}\\[-0.5pt]\textbf{#2}}}

% Pacote: #1 opcoes do corpo, #2 nome do no, #3 e #4 cantos opostos (fit),
% #5 rotulo da aba
\newcommand{\pacote}[5][]{%
  \node[pacotecorpo, fit={#3 #4}, #1] (#2) {};
  \node[pacoteaba, anchor=south west] at (#2.north west) {#5};
}

% No de implantacao (caixa em perspectiva): #1 opcoes do corpo, #2 nome do
% no, #3 e #4 cantos opostos (fit), #5 rotulo (estereotipo + nome)
\newcommand{\nodo}[5][]{%
  \node[nodocorpo, fit={#3 #4}, #1] (#2) {};
  \filldraw[draw=unooNavy, line width=0.85pt, fill=unooNavy!22]
    (#2.north west) -- ++(2.4mm,2.4mm) -- ([shift={(2.4mm,2.4mm)}]#2.north east)
    -- (#2.north east) -- cycle;
  \filldraw[draw=unooNavy, line width=0.85pt, fill=unooNavy!40]
    (#2.north east) -- ++(2.4mm,2.4mm) -- ([shift={(2.4mm,2.4mm)}]#2.south east)
    -- (#2.south east) -- cycle;
  \node[nodolbl, anchor=north west] at (#2.north west) {#5};
}

% Interface fornecida (pirulito): #1 centro da bola
\newcommand{\bola}[1]{%
  \draw[line width=0.85pt, draw=unooInkSoft, fill=unooSurface] #1 circle (1.3mm);}
% Interface requerida (soquete): #1 centro, #2 angulo para onde o soquete ABRE
\newcommand{\soquete}[2]{%
  \draw[line width=0.85pt, draw=unooInkSoft]
    #1 ++(#2+90:2.2mm) arc[start angle=#2+90, end angle=#2+270, radius=2.2mm];}

% Mensagem de diagrama de comunicacao: seta curta ao lado do link.
% #1 posicao do rotulo, #2 ponto medio, #3 angulo da seta, #4 rotulo
\newcommand{\msgcom}[4][above]{%
  \draw[seqmsg] #2 ++(#3+180:4.5mm) -- ++(#3:9mm)
    node[msgcomlbl, midway, #1=1.5pt]{#4};}

\graphicspath{{../../}}

% Rodape com o codigo deste documento, nao o da PU-DOC-001
\renewcommand{\docid}{PU-EST-003}
\renewcommand{\docversao}{1.0}

% --- Selos de status (somente tokens unoo*) ---------------------------------
\newcommand{\selo}[3]{%
  \tikz[baseline=-0.55ex]\node[
    fill=#1, draw=#2, text=#2, line width=0.5pt, rounded corners=2pt,
    inner xsep=4.5pt, inner ysep=1.6pt,
    font=\sffamily\bfseries\fontsize{6.8}{8}\selectfont
  ]{#3};%
}
\newcommand{\implementado}{\selo{unooGreenWash}{unooGreen!75!black}{IMPLEMENTADO}}
\newcommand{\parcial}{\selo{unooOrangeWash}{unooOrange!85!black}{PARCIAL}}
\newcommand{\naoimpl}{\selo{unooRedWash}{unooRed}{NÃO IMPLEMENTADO}}
\newcommand{\naoexiste}{\selo{unooRedWash}{unooRed}{NÃO EXISTE ATUALMENTE}}
\newcommand{\foradorepo}{\selo{unooLavenderWash}{unooNavy}{FORA DO REPOSITÓRIO}}
\newcommand{\variapormissao}{\selo{unooBlueWash}{unooBlue!60!black}{VARIA POR MISSÃO}}

% --- Caixa de resposta de uma pergunta (estilo da historia do mapa) ---------
%  #1 codigo (P01), #2 titulo curto, #3 selo de status
\newtcolorbox{resposta}[3]{
  enhanced, breakable,
  colback=white, colframe=unooGrid,
  borderline west={3pt}{0pt}{unooNavy},
  boxrule=0.5pt, arc=2pt,
  left=10pt, right=10pt, top=6pt, bottom=6pt,
  before skip=10pt, after skip=10pt,
  colbacktitle=unooLavenderWash, coltitle=unooNavy,
  fonttitle=\sffamily\bfseries\small,
  toptitle=4pt, bottomtitle=4pt, lefttitle=10pt,
  title={\texttt{#1}\hspace{6pt}#2\hfill#3}
}

% Campos da resposta: Hoje / Limitações / Planejado / Evidência
\newlist{campos}{description}{1}
\setlist[campos]{
  font=\sffamily\bfseries\footnotesize\color{unooNavy},
  leftmargin=27mm, labelwidth=24mm, labelsep=3mm,
  align=left, itemsep=3pt, topsep=3pt, parsep=1pt
}

% Sugestao de redacao para o TR (requisito observavel + criterio de aceite)
\newtcolorbox{paraotr}[1][]{unoobase, colback=unooBlueWash, colframe=unooBlue,
  title={Proposta de redação para o TR}, #1}

\begin{document}

% ===========================================================================
%  CAPA
% ===========================================================================
\begin{titlepage}
\begin{tikzpicture}[remember picture, overlay]
  \fill[unooNavy] (current page.north west) rectangle ([yshift=-26mm]current page.north east);
  \fill[unooBlue] ([yshift=-26mm]current page.north west) rectangle ([yshift=-28.5mm]current page.north east);
  \fill[unooLavenderWash] (current page.south west) rectangle ([yshift=22mm]current page.south east);
  \draw[unooLavender, line width=0.6pt]
    ([yshift=22mm]current page.south west) -- ([yshift=22mm]current page.south east);
  \node[anchor=west, text=white, font=\sffamily\bfseries\small]
    at ([xshift=22mm, yshift=-13mm]current page.north west)
    {UNOORIA \enspace\textperiodcentered\enspace PRODUTO, ENGENHARIA E LICITAÇÃO};
  \node[anchor=east, text=unooLavenderSoft, font=\sffamily\footnotesize]
    at ([xshift=-22mm, yshift=-13mm]current page.north east)
    {CÓDIGO: PU-EST-003 \enspace|\enspace VERSÃO 1.0};
\end{tikzpicture}

\vspace*{28mm}
\begin{flushleft}
  \includegraphics[width=50mm]{figuras/marca/logo_azul_profundo.png}\\[14mm]
  {\sffamily\bfseries\fontsize{10}{12}\selectfont\color{unooBlue}
    SUBSÍDIO AO TERMO DE REFERÊNCIA}\\[4mm]
  {\sffamily\bfseries\fontsize{25}{30}\selectfont\color{unooNavy}
    Liga Digital}\\[3mm]
  {\sffamily\fontsize{14}{19}\selectfont\color{unooInkSoft}
    Respostas ao questionário de validação para os times: o que a plataforma
    faz hoje, o que falta, o que a V1 promete e a evidência de cada
    afirmação}\\[10mm]
  {\color{unooBlue}\rule{36mm}{2pt}}\\[10mm]
\end{flushleft}

\begin{tcolorbox}[
  enhanced, colback=unooSurface, colframe=unooGrid,
  boxrule=0.5pt, arc=2pt, left=10pt, right=10pt, top=8pt, bottom=8pt
]
\footnotesize\renewcommand{\arraystretch}{1.35}
\begin{tabularx}{\linewidth}{@{}L{40mm}X@{}}
  \textbf{\sffamily\color{unooNavy}Questionário respondido} &
    \emph{Liga Digital --- Questionário de validação para os times},
    versão 1, de 30/10/2026: dez perguntas de Produto e Pedagogia
    (P01--P10), onze de Tecnologia, IA, Segurança e Operação (T01--T11), três
    de Editorial, Formação e Implantação (E01--E03) e os seis compromissos
    obrigatórios da V1 (V1-01 a V1-06).\\
  \textbf{\sffamily\color{unooNavy}Base das respostas} &
    O monorepo da Plataforma Unooria na branch \cd{dev-local}: código,
    testes, suítes e2e, a Documentação Técnica (\cd{PU-DOC-001}), o Mapa de
    Produto (\cd{PU-PRD-001}) e o estudo de carga (\cd{PU-EST-001}).\\
  \textbf{\sffamily\color{unooNavy}Regra de leitura} &
    Cada afirmação aponta a evidência de onde saiu. O que só existe em
    desenvolvimento não é tratado como implementado, e o que o repositório não
    responde (PoC, coleção didática, formação, SLA comercial) está marcado como
    \foradorepo, para o time responsável completar.\\
  \textbf{\sffamily\color{unooNavy}Data de referência} &
    30 de outubro de 2026\\
\end{tabularx}
\end{tcolorbox}

\vfill
\begin{flushleft}
  \footnotesize\sffamily\color{unooMuted}
  Documento complementar à Documentação Técnica da Plataforma Unooria
  (\arq{documentacao-latex/main.pdf}) e ao Mapa de Produto (\cd{PU-PRD-001}).
  Destina-se às equipes que consolidam o Termo de Referência; não é peça do
  edital.
\end{flushleft}
\vspace*{10mm}
\end{titlepage}

\pagenumbering{roman}
\tableofcontents
\cleardoublepage
\pagenumbering{arabic}
\pagestyle{unoomain}

% ===========================================================================
%  CAPITULO 1 — Como ler estas respostas
% ===========================================================================
\chapter{Como ler estas respostas}
\label{cap:como-ler}

Este documento responde, pergunta por pergunta, ao \emph{Questionário de
validação para os times} da Liga Digital (versão 1, 30/10/2026), escrito para
reunir as evidências que faltam ao Termo de Referência (TR) técnico e
pedagógico. As respostas foram levantadas no próprio repositório da
plataforma, e não em apresentações comerciais: o que está aqui como
implementado tem código, teste ou registro de decisão que o sustenta.

\section{Nomes do questionário e nomes da plataforma}

O questionário usa o vocabulário comercial. A plataforma usa outro, e as
evidências citadas estão escritas nele. A correspondência é esta:

\begin{center}
\footnotesize\renewcommand{\arraystretch}{1.3}
\begin{tabularx}{\linewidth}{@{}L{42mm}Y@{}}
\toprule
\textbf{No questionário} & \textbf{Na plataforma e nas evidências}\\
\midrule
Liga Digital & Plataforma Unooria: o hub web (professor, estudante, gestores
  e administração) e as experiências embutidas nele.\\
Missão & Experiência imersiva, identificada por um código como
  \cd{exp08ef1b} (8º ano do EF, 1º bimestre) ou \cd{exp02em1b} (2ª série do
  EM, 1º bimestre). Cada uma simula um aplicativo fictício (Unoovista,
  Unoochat, UnooriaAI, Instanooria\ldots).\\
Ficha da missão & A tela de detalhe da missão no painel do professor,
  alimentada pelo manifesto público da experiência
  (\arq{public/experiencia.json}).\\
Contextualização / personalização / setup & A criação da atividade pelo
  professor: tema e contexto escritos por ele, a partir dos quais a IA gera os
  conteúdos da simulação.\\
Aplicação para a turma & A \emph{liberação} da missão numa turma
  (\adrr{ADR-69}) somada à atividade que o professor põe no ar.\\
Tentativa do estudante & A \emph{rodada} ou sessão do estudante dentro de uma
  atividade.\\
Disponibilidade na rede & O interruptor do painel da Unooria
  (\adrr{ADR-87}): se o produto ou a missão existe naquele contrato, regional
  ou escola.\\
\bottomrule
\end{tabularx}
\end{center}

\section{Estrutura de cada resposta}

Cada pergunta ganha uma caixa com o código do questionário, um selo de status
e cinco campos, na ordem que o questionário pediu:

\begin{campos}
  \item[Hoje] o comportamento que existe, verificado no código.
  \item[Limitações] o que não funciona, só funciona em parte do catálogo, ou
    exige condição operacional.
  \item[Planejado] o que está em desenvolvimento ou decidido, com a fonte.
    Nada aqui é tratado como implementado.
  \item[Evidência] onde conferir: arquivo do repositório, ADR (Capítulo 18 da
    Documentação Técnica, \cd{PU-DOC-001}), história do Mapa de Produto
    (\cd{PU-PRD-001}) ou suíte e2e.
  \item[Lacuna] o que esta resposta não fecha e quem precisa fechar.
\end{campos}

Quando ajuda, a caixa é seguida de uma \textbf{proposta de redação para o TR}:
um requisito escrito como resultado observável com critério de aceite, sem
amarrar arquitetura, marca ou fornecedor, como pedem as observações finais do
questionário.

\section{Selos de status}

\begin{center}
\footnotesize\renewcommand{\arraystretch}{1.4}
\begin{tabularx}{\linewidth}{@{}L{48mm}Y@{}}
\toprule
\textbf{Selo} & \textbf{Significado}\\
\midrule
\implementado & Existe em todo o escopo perguntado, com evidência.\\
\parcial & Existe, mas só em parte do catálogo, sem algum dos controles
  pedidos ou com restrição operacional relevante.\\
\variapormissao & O comportamento difere entre missões; a resposta traz o
  recorte.\\
\naoimpl\newline\naoexiste & Não há código para isso hoje. Quando é
  compromisso da V1, a lacuna está descrita no Capítulo~\ref{cap:v1}.\\
\foradorepo & A resposta depende de material que não está no repositório
  (relatório de PoC, coleção didática, contrato de suporte) e precisa vir do
  time dono.\\
\bottomrule
\end{tabularx}
\end{center}

\begin{atencao}
\textbf{O repositório não responde tudo.} Perguntas sobre a prova de conceito
com o Instituto Paula Souza (P10), a coleção didática e o livro do professor
(E01, E02), a formação e o modelo de implantação (E03) e os compromissos
comerciais de suporte (parte de T09) dependem de documentos que não estão
versionados aqui. Para elas, o documento registra o que a plataforma oferece de
apoio e lista exatamente o que o time responsável precisa anexar.
\end{atencao}

\section{Data e base da verificação}

As respostas refletem a branch \cd{dev-local} em 30 de outubro de 2026. A
plataforma muda toda semana: antes de publicar o TR, confirme em ambiente de
homologação, com dados sintéticos, cada item marcado como implementado, como o
próprio questionário recomenda.

% ===========================================================================
%  CAPITULO 2 — Quadro-resumo
% ===========================================================================
\chapter{Quadro-resumo}
\label{cap:resumo}

Uma linha por item do questionário, com o status e a resposta em uma frase.
O detalhe e as evidências estão na seção indicada.

\begin{center}
\footnotesize\renewcommand{\arraystretch}{1.3}
\begin{tabularx}{\linewidth}{@{}L{11mm}L{36mm}X@{}}
\caption{Status de cada item do questionário}\label{tab:resumo}\\
\toprule
\textbf{Item} & \textbf{Status} & \textbf{Resposta curta}\\
\midrule
\endfirsthead
\toprule
\textbf{Item} & \textbf{Status} & \textbf{Resposta curta}\\
\midrule
\endhead
\multicolumn{3}{@{}l}{\textbf{\sffamily\color{unooNavy}A. Produto e Pedagogia}}\\
P01 & \implementado & Ficha e guia da aula por missão, com objetivos, BNCC,
  duração, roteiro e rubrica; sem campo de bimestre nem de livro.\\
P02 & \variapormissao & Assistente em 4 passos; texto livre e componente
  curricular; \textbf{não aceita anexo}; a IA nunca muda mecânica nem
  rubrica.\\
P03 & \parcial & Atividade salva e reaproveitável na mesma turma;
  \textbf{não duplica nem compartilha} entre turmas ou professores.\\
P04 & \variapormissao & Nada vai à turma sem aprovação; edição, nova geração
  por item e aprovação item a item ou em lote, conforme a missão.\\
P05 & \variapormissao & Liberar, pôr no ar, encerrar e as etapas de cada
  missão; atualização a cada 2 a 5\,s; as etapas só avançam.\\
P06 & \implementado & Missão-base, atividade e rodada; o histórico fica e
  vale a melhor rodada; unicidade no banco evita duplicidade.\\
P07 & \implementado & Avalia ações estruturadas, nunca o texto da conversa;
  IA só classifica intenção em lista fechada, com conjunto de referência.\\
P08 & \naoexiste & Não há recomendação por IA; o painel ordena por regra e o
  agente do professor não lê resultados.\\
P09 & \parcial & Alvo WCAG 2.2 AA; zero violação automática; falta avaliação
  manual, locuções e controle de áudio.\\
P10 & \foradorepo & Nada da PoC no repositório; o time precisa anexar o
  relatório.\\
\midrule
\multicolumn{3}{@{}l}{\textbf{\sffamily\color{unooNavy}B. Tecnologia, IA, Segurança e Operação}}\\
T01 & \parcial & Web sem instalador, responsivo medido; testado só no
  Chromium; sem matriz de compatibilidade.\\
T02 & \parcial & Login próprio com tokens no padrão OIDC; federação WSO2
  preparada e não operacional; sem SAML; revogação espera a sessão vencer.\\
T03 & \implementado & OneRoster 1.2 CSV e planilha, por tela e API;
  idempotente, parcial, com erros por linha; a Secretaria envia, não há
  busca.\\
T04 & \implementado & Seis perfis; autorização na API e isolamento por escola
  no banco; testes de segregação.\\
T05 & \parcial & Regras invioláveis, classificação fechada, aprovação humana;
  sem configuração própria dos filtros do provedor nem registro de
  incidentes.\\
T06 & \parcial & IA recebe texto sem identificação; dados nos EUA com
  salvaguardas da ANPD; \textbf{sem política de retenção}.\\
T07 & \parcial & Auditoria de recursos e de evidência; sem auditoria de login,
  publicação e aprovação de conteúdo.\\
T08 & \parcial & Progresso gravado a cada jogada e retomado; exige conexão
  contínua; sem modo offline.\\
T09 & \parcial & SLOs e alertas definidos; \textbf{teste de carga não
  executado}; sem compromisso de suporte.\\
T10 & \parcial & Relatório por missão e escola; sem filtro de período, turma
  ou bimestre, sem último acesso.\\
T11 & \naoexiste & Professor não consegue jogar a missão; é o V1-02.\\
\midrule
\multicolumn{3}{@{}l}{\textbf{\sffamily\color{unooNavy}C. Editorial, Formação e Implantação}}\\
E01 & \foradorepo & 14 missões em 12 séries; a matriz com o livro é do
  editorial.\\
E02 & \parcial & Ficha, guia imprimível, Ajuda e agente na plataforma; falta
  o recorte do livro.\\
E03 & \foradorepo & A plataforma tem as ferramentas de implantação; o serviço
  não está descrito no repositório.\\
\midrule
\multicolumn{3}{@{}l}{\textbf{\sffamily\color{unooNavy}Compromissos da V1}}\\
V1-01 & \parcial & Faltam filtros, último acesso e baixa participação.\\
V1-02 & \naoimpl & Recurso inteiro a construir, missão por missão.\\
V1-03 & \naoimpl & Entradas prontas; falta a geração e o processo pedagógico.\\
V1-04 & \naoimpl & O tema já está gravado na atividade; depende de V1-03.\\
V1-05 & \implementado & Ajuste do professor prevalece; falta um teste e2e.\\
V1-06 & \parcial & Trilha permanente existe; justificativa opcional e sem
  tela de consulta.\\
\bottomrule
\end{tabularx}
\end{center}

% ===========================================================================
%  CAPITULO 3 — Produto e Pedagogia (P01–P10)
% ===========================================================================
\chapter{Produto e Pedagogia}
\label{cap:produto}

% ---------------------------------------------------------------------------
\section{P01 · O que o professor vê antes de aplicar a missão}
\label{sec:p01}

\begin{resposta}{P01}{Ficha e guia da missão}{\implementado}
\begin{campos}
  \item[Hoje] Duas telas por missão, ambas no painel do professor, seção
    \emph{Entre Redes}, escolhendo ano e unidade:
    \begin{itemize}
      \item \textbf{Ficha da missão} (``Conhecer a missão''), em quatro abas:
        \emph{Como funciona} (os atos da missão), \emph{Antes de rodar}
        (preparo), \emph{O que ela não conta} (opcional) e \emph{Competência e
        formato}. Campos obrigatórios em toda missão publicada: título, resumo,
        unidade, competência, duração, formato, atos e preparo. Os códigos BNCC
        aparecem como etiquetas no cabeçalho (\adrr{ADR-100}, \adrr{ADR-182}).
      \item \textbf{Guia da aula}, em quatro tempos (\emph{Preparar},
        \emph{Conduzir}, \emph{Fechar}, \emph{Consultar}) e versão para
        impressão. Campos obrigatórios: objetivos; habilidades BNCC (código,
        texto, eixo e onde acontece na missão); preparo técnico e pedagógico e
        ``o que não dizer''; roteiro minuto a minuto com o comando de cada
        momento; debate, observação e mediação; síntese, registro,
        desdobramentos e \textbf{rubrica} (critério $\times$ em
        desenvolvimento / adequado / avançado); cuidados, adaptações e dúvidas
        frequentes. Avisos e produto final são opcionais. Um teste automático
        recusa publicar missão sem guia (\adrr{ADR-65}).
    \end{itemize}
    Os critérios por nível de cada habilidade aparecem também no painel de
    habilidades da atividade, nas missões que têm régua de avaliação.
  \item[Limitações] \textbf{Não há campo de bimestre}: a missão é identificada
    por ano/série e \emph{unidade} (1 a 4), e a equivalência unidade =
    bimestre não está formalizada. \textbf{Não há campo de correspondência
    com o livro}: aparece como texto livre em alguns guias (por exemplo,
    ``capítulo 1 do livro, pp. 2--9'' no guia da 2ª série do EM). Materiais
    necessários ficam dentro do preparo, sem campo próprio. Nada disso é
    editável pelo professor: é conteúdo fixo, versionado pela Unooria.
  \item[Evidência] \arq{apps/\B{}hub/\B{}lib/\B{}entre-redes/\B{}catalogo.ts} (ficha);
    \arq{apps/\B{}hub/\B{}lib/\B{}entre-redes/\B{}guias/\B{}} (um guia por missão) e o teste
    \arq{guias/index.test.ts}; suíte \arq{scripts/\B{}e2e/\B{}professor/\B{}08-guias-de-aula.mjs};
    Mapa de Produto, H04.3.2 e H04.4.1.
  \item[Lacuna] Produto: capturas das duas telas de uma missão real para anexo
    (podem ser tiradas da suíte e2e citada). Pedagogia/Editorial: formalizar
    unidade = bimestre e a referência ao livro como campo.
\end{campos}
\end{resposta}

% ---------------------------------------------------------------------------
\section{P02 · O que o professor informa para contextualizar}
\label{sec:p02}

\begin{resposta}{P02}{Formulário de personalização}{\variapormissao}
\begin{campos}
  \item[Hoje] Na ficha, \emph{Nova atividade} abre um assistente de quatro
    passos: \textbf{Turma \seta{} Conteúdo \seta{} Equipes \seta{} Revisar e
    criar} (\adrr{ADR-130}). O passo \emph{Conteúdo} é próprio de cada missão;
    as entradas são texto livre e escolha do componente curricular. Na maioria
    das missões há o atalho \emph{Usar banco padrão}, que carrega conteúdo
    curado pela Unooria; sem saldo de IA ou com o provedor indisponível, a
    tela usa o banco padrão. A Tabela~\ref{tab:p02} resume as entradas e o que
    a IA produz.
  \item[Fixo] Em todas as missões a IA \textbf{não altera} mecânica, etapas,
    gabaritos, rubricas e réguas de habilidade. Exemplos: em \emph{Quem Mandou
    Isso?} a resposta certa de cada mensagem não é editável; no \emph{Porta-Voz
    dos Dados} os números da série e o laudo nem passam pelo modelo; no
    \emph{Enigma da Capa} é o professor quem marca na imagem onde estão os
    objetos.
  \item[Limitações] \textbf{Não é possível anexar documentos} (PDF, plano de
    aula, apostila ou link) para orientar a geração: não existe campo de
    envio de arquivo. Os campos ``nome do arquivo'' que aparecem em duas
    missões são conteúdo fictício da simulação, não anexos. Duas missões não
    têm personalização por IA (\emph{Rastro do Boato} e a versão anterior da
    2ª série do EM).
  \item[Evidência] \arq{apps/\B{}hub/\B{}components/\B{}atividades/\B{}assistente-de-atividade.tsx};
    entradas reais das suítes e2e: ``Boato da merenda cancelada na escola''
    (\arq{estudante/\B{}10-quem-mandou-isso.mjs}), ``Água limpa e saúde no
    bairro'' com quatro tópicos (\arq{estudante/\B{}12-jornal-da-escola.mjs}),
    ``Revolução Industrial na Inglaterra'' (\arq{professor/\B{}02-atividade-do-zero.mjs}).
  \item[Lacuna] O resultado gerado não fica gravado nas suítes (é gerado a
    cada execução). Produto: anexar um par entrada/saída capturado em
    homologação.
\end{campos}
\end{resposta}

\begin{center}
\footnotesize\renewcommand{\arraystretch}{1.3}
\begin{tabularx}{\linewidth}{@{}L{34mm}L{58mm}X@{}}
\caption{Entradas do professor e o que a IA gera, por missão}\label{tab:p02}\\
\toprule
\textbf{Missão} & \textbf{O professor informa} & \textbf{A IA gera}\\
\midrule
\endfirsthead
\toprule
\textbf{Missão} & \textbf{O professor informa} & \textbf{A IA gera}\\
\midrule
\endhead
O Que a Foto Conta (1º EF) & tema fechado: pátio, horta, sala ou parquinho
  & legendas e comentários\\
Enigma da Capa (2º EF) & componente, tema, 3 a 6 palavras-chave
  & a cena e a imagem panorâmica\\
Nina e a Notícia (3º EF) & componente, tema do evento, assunto
  & o texto das 3 notícias (estrutura fixa)\\
Jornal da Escola (4º EF) & componente, tema central, tópicos das 8 notícias
  & 8 notícias e imagens\\
Quem Mandou Isso? (5º EF) & componente, boato fictício, 3 a 6 fatos, contexto
  local (opcional) & 8 mensagens\\
Júri Estadual (6º EF) & componente, título da mostra, 3 a 6 assuntos
  & 40 estandes e imagens\\
A Seleção (7º EF) & evento central, detalhe da cidade (opcional)
  & 12 posts e as falas do podcast\\
Unoovista (8º EF) & título, disciplina, tema, o que já foi trabalhado
  (opcional), proporção & perguntas-chave com nota de equilíbrio; posts e
  imagens\\
Manifesto da IA (8º EF) & dilema da escola (opcional) & textos de apoio\\
Dossiê Invisível (9º EF) & componente, tema, conceitos, contexto (opcional)
  & documentos do dossiê (quantidade fixa)\\
Assistente Invisível (1ª EM) & disciplina, tema, conceitos, foco regional
  (opcional) & as três conversas do assistente\\
Efeito Amplificador (2ª EM) & componente, tema da campanha, 3 a 6 tópicos,
  contexto & conteúdo da campanha\\
Porta-Voz dos Dados (3ª EM) & componente, tema gerador, 3 a 6 objetos de
  conhecimento, bairro (opcional) & textos (números e laudo fixos)\\
Rastro do Boato (6º EF) & --- & sem IA; conteúdo fixo\\
\bottomrule
\end{tabularx}
\end{center}

% ---------------------------------------------------------------------------
\section{P03 · Salvar, reaproveitar e compartilhar}
\label{sec:p03}

\begin{resposta}{P03}{A missão contextualizada como versão reutilizável}{\parcial}
\begin{campos}
  \item[Hoje] A atividade criada fica \textbf{salva} e pertence a
    \textbf{uma turma}. Na mesma turma ela é reaproveitável: a lista
    \emph{Suas atividades} mostra as atividades em andamento, prontas e em
    preparo, e \emph{Usar esta atividade agora} a põe no ar de novo, abrindo
    uma rodada nova e preservando a anterior. Só uma atividade fica no ar por
    turma; a que sai volta a ``pronta'' com tudo guardado. Também é possível
    reaproveitar as equipes da atividade anterior.
  \item[Quem faz o quê] Vê e conduz: o professor que leciona na turma (outra
    turma responde como inexistente). Pela interface, só o professor edita;
    a coordenação lê. Exclusão só existe numa missão (Unoovista), com
    confirmação pela digitação do título. Não há fluxo de aprovação por
    terceiros. O conteúdo aprovado fica gravado na própria atividade; se a
    Unooria atualizar a missão-base, a atividade já criada mantém o conteúdo
    aprovado.
  \item[Limitações] \textbf{Duplicar para outra turma, outro período ou outro
    professor não existe atualmente} (exceção: no Unoovista a atividade pode
    ser limpa e usada com outra turma). \textbf{Compartilhar entre
    professores, escola, rede ou biblioteca comum não existe atualmente.} No
    exemplo do questionário, a professora refaz a personalização para o
    8º ano B; o banco padrão da Unooria é o único conteúdo comum.
  \item[Planejado] Nada registrado no Mapa de Produto nem no \emph{backlog}.
  \item[Evidência] Mapa de Produto, H06.13.1 e H06.13.3; \adrr{ADR-130},
    \adrr{ADR-146}, \adrr{ADR-155}.
\end{campos}
\end{resposta}

% ---------------------------------------------------------------------------
\section{P04 · Curadoria do conteúdo gerado}
\label{sec:p04}

\begin{resposta}{P04}{Revisão antes de a turma ver}{\variapormissao}
\begin{campos}
  \item[Hoje] Regra comum: \textbf{nada gerado chega à turma sem aprovação do
    professor}; sem ela, a plataforma recusa pôr a atividade no ar. O
    professor vê cada item gerado. Os controles variam:
    \begin{itemize}
      \item \textbf{Editar e regenerar item a item}, com aprovação final do
        pacote inteiro: Assistente Invisível, Porta-Voz dos Dados e Dossiê
        Invisível (que marca ``modificado pelo docente''); na Nina só a
        manchete é editável.
      \item \textbf{Duas fases}, textos e depois imagens (\adrr{ADR-146},
        \adrr{ADR-147}): no Jornal da Escola cada imagem é aprovada
        individualmente e ``gerar outra imagem'' a reabre; na Seleção a troca é
        por post e a aprovação das imagens, em lote.
      \item \textbf{Aprovação em lote}: Júri Estadual e imagens da Nina; depois
        de aprovada, a imagem não muda.
      \item \textbf{Unoovista}: filtrar, regerar texto ou imagem, excluir post;
        ativar é recusado se faltar conteúdo.
      \item As falas do podcast são revisadas obrigatoriamente na criação
        (\adrr{ADR-150}, \adrr{ADR-204}).
    \end{itemize}
    Rejeitar um item é pedir nova geração manual daquele item; não há nova
    geração automática. \textbf{Registro}: a maioria das missões grava quem
    aprovou e quando, e o conteúdo aprovado fica salvo na atividade.
    \textbf{Verificação adicional}: regras automáticas conferem estrutura e
    critérios pedagógicos da saída antes de mostrá-la; na Seleção, uma
    segunda leitura por IA confere se o post editado manteve o enquadramento.
  \item[Limitações] Três missões não guardam o autor da aprovação (Nina guarda
    só a data; Manifesto, um indicador sim/não; Unoovista, o dono e o
    estado). Edições não são versionadas. Não há moderação dedicada nem
    configuração própria dos filtros do provedor (T05). O Mapa de Produto
    registra que as imagens geradas precisam de revisão de conteúdo antes da
    primeira rede (H06.12.7).
  \item[Evidência] \arq{apps/\B{}hub/\B{}app/\B{}(plataforma)/\B{}professor/\B{}entre-redes/\B{}AGENTS.md};
    \cd{PU-DOC-001}, Cap. 11; suítes e2e de estudante 10 a 19 e
    \arq{professor/\B{}02-atividade-do-zero.mjs}.
\end{campos}
\end{resposta}

% ---------------------------------------------------------------------------
\section{P05 · Comandos do professor durante a aula}
\label{sec:p05}

\begin{resposta}{P05}{Condução em sala}{\variapormissao}
\begin{campos}
  \item[Hoje] Em toda missão há três comandos de base: \textbf{liberar} a
    missão na turma (sem isso, o estudante vê o cartão com cadeado, e
    digitar o endereço também não abre: \adrr{ADR-69}), \textbf{pôr a
    atividade no ar} e \textbf{encerrar} (fecha a entrada sem expulsar quem
    está jogando, e converte as marcas pendentes em nível). Depois, cada
    missão tem as etapas próprias, conduzidas pelo professor
    (Tabela~\ref{tab:p05}). O estudante vê a tela de espera, a tela
    congelada (pausa) ou a etapa nova; quem está atrasado continua no próprio
    caminho. O painel da aula mostra ``ainda não entrou'', ``concluíram
    $x$/$y$'', a fila por passo e o tempo parado; com equipes, conta por
    computador.
  \item[Limitações] A atualização é por consulta periódica, não em tempo
    real: o estudante relê a etapa a cada 3 a 5 segundos e o painel do
    professor a cada 2 a 5 segundos (\adrr{ADR-93}); todos precisam estar
    conectados. \textbf{As etapas só avançam}: voltar ou pular é recusado, e
    repetir o comando não muda nada. Programar abertura e fechamento ainda não
    existe (H04.2.3, a fazer).
  \item[Evidência] Uma suíte e2e por missão em \arq{scripts/e2e/estudante/};
    \arq{scripts/\B{}e2e/\B{}professor/\B{}06-resultados-das-missoes.mjs}; Mapa de
    Produto, H04.2.1 e H04.2.2.
\end{campos}
\end{resposta}

\begin{center}
\footnotesize\renewcommand{\arraystretch}{1.3}
\begin{tabularx}{\linewidth}{@{}L{42mm}X@{}}
\caption{Comandos de condução, por missão}\label{tab:p05}\\
\toprule
\textbf{Missão} & \textbf{Etapas e comandos do professor}\\
\midrule
\endfirsthead
\toprule
\textbf{Missão} & \textbf{Etapas e comandos do professor}\\
\midrule
\endhead
O Que a Foto Conta & configurar, aprovar o pacote, destravar etapa, reiniciar\\
Enigma da Capa & liberar a virada\\
Nina e a Notícia & aposta \seta{} filtro \seta{} inspeção; pausar e retomar a turma\\
Jornal da Escola & exploração \seta{} resultado \seta{} revisão\\
Quem Mandou Isso? & revelação, síntese (opcional), reflexão; projetar o termômetro da turma\\
Júri Estadual & exploração \seta{} galeria (liberar auditoria)\\
Rastro do Boato & liberar a aula 2; pausar a simulação\\
A Seleção & exploração \seta{} mural \seta{} fichas \seta{} explicação\\
Manifesto da IA & pausar para debate; retomar\\
Dossiê Invisível & exploração \seta{} revelação \seta{} revoto \seta{} explicação (sem pausa)\\
Assistente Invisível & testes \seta{} simulação \seta{} explicar \seta{} publicar, com duas pausas
  de debate opcionais; a simulação exige todas as apostas feitas\\
Efeito Amplificador & exploração \seta{} pausa para debate \seta{} explicar\\
Porta-Voz dos Dados & exploração \seta{} pausa \seta{} liberação da errata\\
\bottomrule
\end{tabularx}
\end{center}

% ---------------------------------------------------------------------------
\section{P06 · Nova aplicação e novas tentativas}
\label{sec:p06}

\begin{resposta}{P06}{Missão-base, atividade e rodada}{\implementado}
\begin{campos}
  \item[Hoje] Os três objetos existem: a \textbf{missão-base} (o esqueleto
    publicado pela Unooria), a \textbf{atividade} (a aplicação
    contextualizada, de uma turma) e a \textbf{rodada} (cada vez que a
    atividade vai ao ar). (1) O professor pode criar \textbf{várias atividades
    da mesma missão para a mesma turma}; só uma fica no ar por vez.
    (2) Pôr no ar de novo uma atividade que já rodou abre a \textbf{rodada
    seguinte}, e o estudante refaz a missão; a rodada anterior fica guardada
    e toda rodada mede tudo de novo (\adrr{ADR-145}). \emph{Zerar a turma}
    (\adrr{ADR-155}) abre nova rodada e apaga as marcas e os níveis daquela
    atividade, inclusive ajustes do professor, para recomeçar limpo (útil
    após erro); o histórico de eventos permanece. (3) \textbf{O que vale}: na
    mesma rodada, o melhor registro; na atividade, a melhor rodada por
    habilidade; no painel de habilidades, o melhor resultado entre todas as
    atividades e rodadas da missão (\adrr{ADR-144}). Ajuste do professor
    prevalece sobre registro próprio, que prevalece sobre o da equipe.
    (4) \textbf{Duplicidade}: o banco aceita um único resultado por escola,
    missão, atividade, rodada, estudante e habilidade; liberar e pôr no ar são
    idempotentes; atividade e rodada nunca vêm do navegador, são definidas no
    servidor.
  \item[Limitações] Exceções: no Unoovista vale o primeiro envio; em
    múltipla escolha vale a primeira resposta, salvo régua que diga o
    contrário. Não há regra específica para ano letivo seguinte: a atividade
    é da turma, e a turma nova é outra turma (P03).
  \item[Evidência] Mapa de Produto, H14.2.2 e H14.2.4; testes
    \arq{apps/\B{}api/\B{}tests/\B{}evidencias-consolidado-postgres.spec.ts}; suítes
    \arq{professor/\B{}13-painel-de-habilidades.mjs} e
    \arq{professor/\B{}11-perfil-do-aluno.mjs}.
\end{campos}
\end{resposta}

% ---------------------------------------------------------------------------
\section{P07 · Evidência quando o estudante conversa com a IA}
\label{sec:p07}

\begin{resposta}{P07}{Conversa livre e avaliação}{\implementado}
\begin{campos}
  \item[Hoje] Três missões têm texto livre do estudante dirigido à IA: o
    \emph{Efeito Amplificador} (2ª EM), o \emph{Dossiê Invisível} (9º EF) e a
    versão anterior da 2ª série (fora do catálogo, mantida para turmas que já
    a usavam). O Assistente Invisível (1ª EM) \textbf{não} tem conversa
    livre: o estudante lê conversas preparadas. (1) \textbf{A avaliação usa
    eventos estruturados, nunca o conteúdo da conversa}: o que pontua são
    ações (publicar, avisar, proteger, escolher um caminho). Nas missões
    atuais, a IA só traduz a mensagem numa intenção de lista fechada, que
    decide o caminho da história; ``a IA nunca decide o que a régua lê''
    (\adrr{ADR-192}). (2) Os critérios estão nas \textbf{réguas de
    avaliação}, uma por missão, escritas a partir da régua pedagógica, com
    códigos BNCC conferidos automaticamente contra o currículo e versão
    gravada em cada marca. (3) Há LLM classificando em dois pontos: a
    intenção da mensagem (temperatura 0, saída em formato fechado e validado,
    confiança mínima de 0,7, árvore de regras quando a IA falha) e, na
    \emph{Seleção}, a legenda escrita pelo estudante, lida em três
    indicadores; o nível é decidido por regra no servidor e, se a IA falhar,
    fica ``não avaliado'' para o professor ajustar. A consistência é medida
    por conjunto de referência com a IA real: 88/89 e 89/89 acertos, zero
    crédito indevido. (4) O texto do estudante não vai a painel, telemetria
    nem registro técnico; na versão anterior da 2ª série o professor lê a
    transcrição de cada estudante, e o gestor só vê o agregado.
  \item[Exemplo] Do conjunto de referência: a mensagem ``essa frase da empatia
    é sua mesmo?'' é classificada como a intenção ``auditar autoria'', que
    leva ao caminho em que o estudante investiga a origem do texto; a
    evidência vem desse caminho e das ações seguintes, não da frase.
  \item[Limitações] O classificador não é determinístico: um erro custa ao
    estudante reformular a mensagem, nunca um crédito indevido.
  \item[Evidência] \arq{apps/\B{}api/\B{}src/\B{}efeito-amplificador/\B{}AGENTS.md};
    \arq{apps/\B{}api/\B{}scripts/\B{}medir-classificador-exp02em1b.ts};
    \arq{apps/\B{}api/\B{}src/\B{}evidencias/\B{}fichas/\B{}}; \arq{tests/\B{}evidencias-fichas.spec.ts};
    \adrr{ADR-172}, \adrr{ADR-191}--\adrr{ADR-193}.
\end{campos}
\end{resposta}

% ---------------------------------------------------------------------------
\section{P08 · Qualidade das recomendações pedagógicas}
\label{sec:p08}

\begin{resposta}{P08}{Recomendações geradas por IA}{\naoexiste}
\begin{campos}
  \item[Hoje] \textbf{Recomendações de intervenção geradas por IA a partir das
    dificuldades da turma não existem atualmente.} O que existe: o
    \textbf{painel de habilidades} mostra, por habilidade, quantos estudantes
    atingiram, estão em processo, não atingiram ou não foram avaliados, e
    ordena por regra fixa (sem IA) para o professor ver onde reforçar; numa
    missão a IA \textbf{sugere um nível} para a reflexão do estudante,
    rotulado ``sugestão da IA'' até o professor confirmar (\adrr{ADR-167});
    o \textbf{agente de IA do professor} responde sobre BNCC e sobre o acervo
    de livros, citando código e página, mas não acessa resultados da turma
    nem evidências.
  \item[Prompts] Ficam no repositório, só no servidor, versionados pelo
    histórico do código; apenas uma missão grava a versão do prompt junto do
    resultado. A autoria pedagógica dos prompts não está registrada nos
    arquivos.
  \item[Observado $\times$ sugerido] Já há separação explícita: sugestão da IA
    é rotulada e só vale após confirmação; ``não avaliado'' nunca conta como
    ``não atingiu''; o agente só cita habilidade que a busca devolveu e avisa
    quando o código não existe.
  \item[Limitações] O agente do professor está parcial e ainda não está no ar
    nos ambientes da rede.
  \item[Evidência] \arq{apps/\B{}hub/\B{}lib/\B{}painel-de-habilidades.ts};
    \arq{apps/\B{}agente-professor/\B{}agente\_professor/\B{}instrucoes.py} e o conjunto de
    avaliação \arq{avaliacao/\B{}agente\_professor.evalset.json} (16 casos);
    Mapa de Produto, F08.4 e H14.4.5.
  \item[Lacuna] É o compromisso V1-03/V1-04 (Capítulo~\ref{cap:v1}).
\end{campos}
\end{resposta}

% ---------------------------------------------------------------------------
\section{P09 · Acessibilidade e inclusão}
\label{sec:p09}

\begin{resposta}{P09}{Recursos por tipo de missão}{\parcial}
\begin{campos}
  \item[Hoje] Alvo declarado: \textbf{WCAG 2.2 AA} (\adrr{ADR-162}), com
    verificação automática (axe-core) nas regras A e AA das versões 2.0, 2.1
    e 2.2, medindo a página e cada missão embutida. Primeira medição: 65 telas
    do hub e a abertura de 11 missões; as violações encontradas foram
    corrigidas no design system (contraste de texto, foco visível, alvo de
    toque, relações de acessibilidade das abas), e a última medição
    registrada é de \textbf{zero violação} nessas telas. Recursos presentes:
    \textbf{texto alternativo} nas imagens geradas por IA, derivado da
    descrição que gerou a imagem, nunca vazio (\adrr{ADR-163}), e nas fotos
    fixas; \textbf{transcrição integral sempre visível} nos áudios de
    podcast, pensada também para sala sem fone; anúncio de mudança para
    leitor de tela em várias missões; respeito à preferência de movimento
    reduzido em várias missões. \textbf{1º ano, não alfabetizados}: instrução
    em caixa alta com até 8 palavras e redundância tripla (texto, pictograma e
    locução), botão grande de repetir o áudio, uma locução por vez.
  \item[Limitações] Só houve verificação automática: \textbf{falta a passada
    manual com teclado e leitor de tela}, que é a que autoriza declarar AA
    (risco R-07). A Vila do Lume (jogo 3D) está fora da medição. As
    \textbf{locuções ainda não foram gravadas} (a missão do 1º ano cai para
    texto e ícone; Manifesto e Rastro do Boato aguardam áudio). \textbf{Não
    há controle para ligar e desligar áudio}: ouvir é sempre opcional (o
    áudio só toca por comando) e o texto sempre o substitui. Não há legenda de
    vídeo nem ampliação de fonte própria (vale o zoom do navegador). A
    preferência de movimento reduzido não está em todas as missões.
  \item[Planejado] Passada manual de teclado e leitor de tela por papel;
    tornar a verificação bloqueante no catálogo de componentes; gravação das
    locuções (responsável: Conteúdo).
  \item[Evidência] \arq{scripts/\B{}e2e/\B{}lib/\B{}acessibilidade.mjs};
    \arq{scripts/\B{}e2e/\B{}ferramentas/\B{}medir-acessibilidade.mjs};
    \cd{PU-DOC-001}, Cap. 16 (estado ``parcial'') e Cap. 19 (R-07).
  \item[Lacuna] O relatório da última medição não é guardado no repositório;
    gerar e anexar um novo antes do TR.
\end{campos}
\end{resposta}

\begin{paraotr}
A solução deve atender à WCAG 2.2 nível AA (ou ABNT NBR 17225) nas telas de
estudante, professor e gestor, com alternativa textual a todo áudio e imagem.
Aceite: relatório de verificação automática sem violações e relatório de
avaliação manual com teclado e leitor de tela, com plano de correção datado
para as barreiras encontradas.
\end{paraotr}

% ---------------------------------------------------------------------------
\section{P10 · Resultados da PoC com o Instituto Paula Souza}
\label{sec:p10}

\begin{resposta}{P10}{PoC}{\foradorepo}
\begin{campos}
  \item[Hoje] O repositório não contém relatório, apresentação, instrumento de
    pesquisa nem dados da PoC; não há menção ao Instituto Paula Souza no
    código nem na documentação. Também não há instrumentação de pesquisa
    (questionário de satisfação, medida de aprendizagem pré e pós) na
    plataforma.
  \item[O que a plataforma oferece] Para uma futura avaliação, os dados que a
    plataforma já registra podem servir de medida objetiva de conclusão e de
    desempenho por habilidade: progresso e conclusão por estudante e missão,
    evidências por habilidade e telemetria das missões (T10).
  \item[Lacuna] Time de Produto/Pedagogia: anexar período, objetivos,
    missões testadas, número e perfil de participantes, metodologia,
    instrumentos, resultados, limitações e mudanças de produto. Sem esse
    material, o TR não deve declarar ganho de aprendizagem, engajamento ou
    aceitação.
\end{campos}
\end{resposta}

% ===========================================================================
%  CAPITULO 4 — Tecnologia, IA, Segurança e Operação (T01–T11)
% ===========================================================================
\chapter{Tecnologia, IA, Segurança e Operação}
\label{cap:tecnologia}

% ---------------------------------------------------------------------------
\section{T01 · Navegador, dispositivo e conectividade}
\label{sec:t01}

\begin{resposta}{T01}{Combinações suportadas}{\parcial}
\begin{campos}
  \item[Hoje] A plataforma é 100\% web: o hub é uma aplicação Next.js e cada
    missão é um pacote estático servido por CDN, sem instalador. As suítes
    automatizadas de ponta a ponta rodam no \textbf{Chromium} (Playwright
    1.62), em 1280$\times$950 (estudante) e 1440$\times$1000 (professor). A
    responsividade do hub foi medida nos seis papéis em \textbf{375, 768 e
    1280\,px}, em cerca de 80 rotas, e as telas que rolavam de lado foram
    corrigidas (\adrr{ADR-159}). Missões desenhadas para celular declaram
    \cd{viewport} no manifesto e rodam numa moldura de 375$\times$812 dentro do
    hub. Os arquivos das missões têm caminho imutável por versão e ficam em
    cache na CDN; só a página de entrada expira em 60\,s.
  \item[Limitações] Não há matriz formal de compatibilidade: Firefox, Safari
    (WebKit) e Edge não são testados automaticamente, e não há versão mínima
    declarada nem sistema operacional homologado. Não há orçamento de tamanho
    de pacote nem teste sob banda reduzida. A medição de responsividade não
    cobriu modais nem telas com dados reais de aula, e o script dela não está
    versionado. Algumas missões foram desenhadas para computador (o Jornal da
    Escola, por exemplo, não declara \cd{viewport}). A Vila do Lume (jogo 3D)
    exige WebGL, teclado e mouse e não tem controle por toque.
  \item[Planejado] Nada registrado além do gatilho da \adrr{ADR-159} (tela nova
    que role de lado em 375\,px).
  \item[Evidência] \adrr{ADR-159}; \arq{scripts/e2e/lib/roteiro.mjs};
    \arq{scripts/\B{}e2e/\B{}estudante/\B{}16-assistente-invisivel.mjs} (emulação de
    celular 412$\times$915); RNF-03 em \cd{PU-DOC-001}, Cap. 3.
  \item[Lacuna] Engenharia precisa executar e anexar a matriz (Chrome, Edge,
    Firefox e Safari recentes; Android e iOS; Chromebook) antes de o TR fixar
    versões mínimas.
\end{campos}
\end{resposta}

\begin{paraotr}
A solução deve funcionar, sem instalação de software, nas duas últimas versões
estáveis de Chrome, Edge, Firefox e Safari, em telas a partir de 360\,px de
largura para as funções do estudante e a partir de 1280\,px para a condução da
aula. Aceite: roteiro de homologação executado pela contratada em cada
combinação, com registro de tela e lista de restrições por missão.
\end{paraotr}

% ---------------------------------------------------------------------------
\section{T02 · Autenticação e identidade da rede}
\label{sec:t02}

\begin{resposta}{T02}{Login, vínculo e revogação}{\parcial}
\begin{campos}
  \item[Hoje] O login em produção é por \textbf{usuário e senha próprios} da
    plataforma: senha guardada com \emph{hash} bcrypt, limite de 8 tentativas
    por minuto por usuário, mensagem de recusa genérica (\adrr{ADR-97}) e troca
    obrigatória da senha provisória. Depois do login, a plataforma emite
    \emph{tokens} assinados (RS256, validade de 15 minutos) com usuário,
    escola, turma principal e papel, e publica a chave pública num endereço
    JWKS padrão; a API valida emissor, audiência e assinatura. Esse desenho
    segue o padrão \textbf{OpenID Connect}, de modo que trocar o emissor por um
    provedor de identidade da rede é troca de configuração. Escola, turma e
    papel vêm das lotações e matrículas cadastradas pela importação (T03);
    quem trabalha em mais de uma escola escolhe a escola ativa no cabeçalho,
    validada contra as lotações. A sessão expira após 2 horas sem uso.
  \item[Limitações] A federação com o provedor da rede (WSO2 Identity Server,
    via OIDC) está \textbf{preparada mas não operacional}: o conector existe e
    liga por configuração, mas o mapeamento dos atributos da rede ainda não
    foi feito. \textbf{SAML não existe atualmente.} Desativar uma conta
    impede novos logins na hora, mas uma sessão já aberta continua válida até
    expirar (até 2 horas de inatividade); não há revogação ativa nem
    desconexão global. Perder uma turma vale na hora (cada pedido confere a
    matrícula); perder a escola ou mudar de papel vale no próximo login.
  \item[Planejado] Ligar o WSO2 por contrato (\arq{docs/INFRA-IDENTITY.md}:
    um \emph{tenant} por contrato, emissor e JWKS próprios). Decisões
    abertas: base de usuários mestre ou federação com criação na primeira
    entrada (Google Workspace, Gov.br), formato do usuário e tela de login.
  \item[Evidência] \cd{PU-DOC-001}, Cap. 5, diagramas \emph{Três relógios:
    token, renovação e sessão} e \emph{A jornada do estudante, do login à
    missão}; \arq{apps/\B{}api/\B{}tests/\B{}senha-provisoria.spec.ts};
    \arq{scripts/\B{}e2e/\B{}publico/\B{}02-login-e-papeis.mjs}; \adrr{ADR-95},
    \adrr{ADR-96}.
  \item[Lacuna] Decidir se encerrar lotação ou desativar conta deve derrubar a
    sessão aberta (hoje não derruba).
\end{campos}
\end{resposta}

\begin{paraotr}
A solução deve oferecer autenticação própria e federação com o provedor de
identidade da contratante por OpenID Connect; SAML 2.0, quando exigido,
deve ser declarado como item opcional. A desativação de um usuário deve
impedir o acesso a dados em até \textit{[prazo a definir]} minutos, inclusive
em sessão já aberta. Aceite: demonstração do login federado com usuário de
teste da rede e da revogação dentro do prazo.
\end{paraotr}

% ---------------------------------------------------------------------------
\section{T03 · Cadastro de escolas, turmas e pessoas}
\label{sec:t03}

\begin{resposta}{T03}{OneRoster e importação}{\implementado}
\begin{campos}
  \item[Hoje] \textbf{OneRoster 1.2, formato CSV}: os sete arquivos do padrão
    (\cd{manifest}, \cd{orgs}, \cd{academicSessions}, \cd{classes},
    \cd{users}, \cd{enrollments}, \cd{roles}), um por envio, em modo
    \cd{bulk} ou \cd{delta}, pela tela do gestor ou por chamada autenticada à
    API. Aceita a extensão brasileira de CPF. Há também uma planilha
    simplificada (matrícula, nome, usuário, e-mail, papel, turma, série,
    turno; até 50 mil linhas). Modelos e exemplos de cada arquivo são
    baixáveis pela própria plataforma. Regras de carga: reimportar o mesmo
    lote não muda nada (idempotente, pela chave da rede e pelo CPF); troca de
    turma encerra a matrícula anterior sem apagá-la; CPF já existente em outra
    escola vira transferência (estudante) ou lotação adicional (professor);
    \cd{tobedeleted} desativa a pessoa sem apagar; uma carga completa
    \textbf{não} desativa quem ficou de fora, só os lista como candidatos;
    CPF repetido ou inválido é erro de linha; linhas válidas entram e as
    inválidas voltam com o motivo; um arquivo fora de ordem informa qual
    arquivo falta. O retorno traz contadores de criados, atualizados,
    transferidos, desativados e erros.
  \item[Limitações] \textbf{A plataforma não busca mudanças}: é a Secretaria
    (ou o gestor) quem envia as atualizações. Não existe a API REST do
    OneRoster nem sincronização agendada, e o pacote zip do padrão não é
    aceito (um arquivo por envio). Limites: 20 envios por minuto (OneRoster)
    e 5 por minuto (planilha). A senha do arquivo é ignorada.
  \item[Planejado] Nada registrado.
  \item[Evidência] \arq{apps/\B{}api/\B{}tests/\B{}oneroster-import.spec.ts} (cerca de 25
    casos, entre eles ``carga completa incompleta não desativa ninguém'' e
    ``reaplicar o mesmo lote não muda nada''),
    \arq{rostering-import.spec.ts}, \arq{rede-cpf-transferencia.spec.ts};
    suíte \arq{scripts/\B{}e2e/\B{}gestor\_contrato/\B{}02-oneroster.mjs};
    \arq{apps/\B{}api/\B{}src/\B{}rostering/\B{}AGENTS.md}.
\end{campos}
\end{resposta}

\begin{paraotr}
A contratada deve aceitar a carga de escolas, turmas, pessoas e vínculos no
padrão OneRoster 1.2 (CSV), por interface e por API autenticada, com
importação parcial, relatório de erros por linha e carga idempotente. A
atualização é de iniciativa da contratante, salvo integração adicional
pactuada. Aceite: carga de teste sintética com inclusão, troca de turma,
transferência, desligamento e linhas inválidas, conferida contra o relatório.
\end{paraotr}

% ---------------------------------------------------------------------------
\section{T04 · Cada usuário vê só o que lhe compete}
\label{sec:t04}

\begin{resposta}{T04}{Perfis e segregação de dados}{\implementado}
\begin{campos}
  \item[Hoje] Seis perfis: \textbf{estudante} (a própria matrícula, nunca
    gabarito nem dados de colega); \textbf{professor} (só as turmas em que
    leciona); \textbf{gestor da escola} (a própria escola);
    \textbf{gestor regional} (as escolas da regional, uma por vez, e o
    panorama somado); \textbf{gestor do contrato} (a rede inteira);
    \textbf{administração Unooria} (o anterior, mais disponibilidade de
    recursos, quotas e custos de IA). Ninguém cria o próprio nível nem um
    acima (\adrr{ADR-80}). A autorização é aplicada \textbf{na API, em toda
    chamada}, em três travas: o papel exigido pela rota; o isolamento por
    escola no próprio banco (\emph{Row-Level Security}), com o escopo tirado
    do \emph{token} assinado e não do parâmetro; e, para o professor, a
    conferência da matrícula na turma a cada pedido. Um identificador de
    outra escola manipulado na chamada cai no isolamento do banco ou volta
    403. Exportações e painéis usam as mesmas rotas protegidas. Cada contrato
    tem a própria infraestrutura (isolamento físico entre redes). Quem
    trabalha em várias escolas acumula lotações e troca a escola ativa;
    transferência de estudante encerra lotação e matrícula de origem, e o
    professor de origem deixa de vê-lo.
  \item[Limitações] Três tabelas de cadastro (pessoas, turmas e lotações)
    ficam fora do isolamento do banco: nelas, a trava é só o filtro por escola
    na consulta, e uma falha dessas já foi encontrada e corrigida
    (\adrr{ADR-49}). A revogação de sessão tem o atraso descrito em T02.
  \item[Evidência] Matriz papel $\times$ rota $\times$ tela em \cd{PU-DOC-001},
    Cap. 7 (``Nível 3''); testes \arq{tenancy.spec.ts} (``escola A não alcança
    linha da escola B nem por chave direta''), \arq{matriculas-tenancy.spec.ts},
    \arq{professor-por-matricula-tenancy.spec.ts}, \arq{planos-tenancy.spec.ts},
    \arq{rede-cpf-transferencia.spec.ts}, além de 32 especificações que
    provam respostas 403; suítes e2e de gestor de escola (quem trabalha em
    duas escolas) e de contrato (transferência e perfil portátil).
\end{campos}
\end{resposta}

\begin{paraotr}
A solução deve aplicar a autorização no servidor, em toda requisição,
impedindo que um usuário acesse dados de escola, turma ou estudante fora do
seu vínculo, inclusive por manipulação de identificadores. Aceite: roteiro de
teste de segregação executado pela contratada com dois usuários de escolas
diferentes, incluindo chamadas diretas à API com identificadores alheios, e
relatório de teste de intrusão independente antes da produção.
\end{paraotr}

% ---------------------------------------------------------------------------
\section{T05 · Conteúdo inadequado nas duas modalidades de IA}
\label{sec:t05}

\begin{resposta}{T05}{Guardrails}{\parcial}
\begin{campos}
  \item[Hoje] Toda chamada de IA passa por uma camada única no servidor;
    nenhuma instrução de sistema chega ao navegador. \textbf{Geração que o
    professor revisa}: toda instrução de geração carrega um bloco de regras
    invioláveis (pessoas fictícias; sem sexo, violência gráfica, ódio,
    drogas, automutilação ou palavrões; sem estudante ou escola reais); o
    resultado passa por validação automática com listas de termos e perfis
    proibidos; o texto do professor entra como dado delimitado e com tamanho
    máximo, uma barreira contra instruções injetadas; e o professor aprova
    antes de a turma ver (P04). \textbf{Estudante falando com a IA}: nas
    missões em que há texto livre, a IA \emph{classifica} a mensagem numa
    lista fechada de intenções e a fala devolvida vem de um catálogo escrito
    pela equipe; com confiança abaixo de 0,7, lentidão, falta de saldo ou
    recusa, vale uma árvore de regras determinística. Pedidos como ``ignore as
    instruções'' são detectados e tratados como fora do tema sem chegar à IA.
    A série é tirada da turma do estudante, nunca do pedido. Há quota diária
    por perfil e limite de chamadas. \textbf{Temas sensíveis}: o check-in de
    bem-estar é o caso mais cercado (itens fixos, bloqueio de vocabulário
    clínico e de links, sinal de risco pula a IA e mostra texto fixo com CVV
    188 e SAMU 192). \textbf{Dados pessoais}: na missão Manifesto, CPF,
    telefone, e-mail, CEP, cartão e @perfil são detectados e o texto é
    recusado sem ser gravado. Golpes, \emph{fake news} e \emph{cyberbullying}
    são tratados como objeto de estudo dentro do roteiro da missão.
  \item[Limitações] \textbf{Não há configuração própria dos filtros de
    segurança do provedor}: valem os padrões do Gemini, e a recusa por
    segurança é só contada nas métricas, sem tratamento específico. Não há
    moderação dedicada de entrada e saída. O tutor genérico devolve o texto do
    modelo sem validação de saída. O filtro de dados pessoais existe só numa
    missão e não detecta nomes próprios. \textbf{Não há registro formal de
    incidentes de IA.}
  \item[Provedor $\times$ Unooria] Do provedor: os filtros padrão do modelo. Da
    Unooria: instruções de sistema, regras invioláveis, classificação fechada,
    validação da saída, detecção de dados pessoais, quotas e a aprovação
    humana.
  \item[Evidência] \arq{apps/\B{}api/\B{}src/\B{}geracao/\B{}prompts/\B{}guardrails.ts};
    \arq{apps/\B{}api/\B{}src/\B{}manifesto-ia/\B{}regras/\B{}moderacao.ts};
    \adrr{ADR-191}--\adrr{ADR-193}; \cd{PU-DOC-001}, Cap. 11 (guardrails do
    check-in); testes \arq{ai-tutor-prompt.spec.ts},
    \arq{manifesto-ia-regras.spec.ts}, \arq{bemestar-regras.spec.ts}; conjunto
    de referência do classificador da 2ª série (88/89 e 89/89 acertos, zero
    crédito indevido).
  \item[Lacuna] Configurar e versionar os filtros do provedor; estender a
    detecção de dados pessoais a toda entrada livre de estudante; criar
    registro e rito de incidentes.
\end{campos}
\end{resposta}

% ---------------------------------------------------------------------------
\section{T06 · Dados pessoais e fluxos para a IA}
\label{sec:t06}

\begin{resposta}{T06}{Mapa de dados}{\parcial}
\begin{campos}
  \item[Hoje] A rede é \textbf{controladora} e a Unooria \textbf{operadora};
    a base legal é a execução de política pública, não o consentimento
    (página \rota{/privacidade}, \adrr{ADR-70}). Categorias tratadas:
    cadastro (nome, usuário, e-mail, CPF, matrícula, escola, turma),
    atividade nas missões, evidências por habilidade, textos escritos nas
    missões e, em missões específicas, conversa com a IA. Infraestrutura:
    Google Cloud em \textbf{\cd{us-east1}} (EUA), com banco gerenciado de alta
    disponibilidade e recuperação a um ponto no tempo; IA no Vertex AI
    (endpoint global) e voz ao vivo da Ajuda em \cd{us-central1}; avatar da
    Ajuda (LiveAvatar/HeyGen, EUA) só para profissionais. A transferência
    internacional é declarada na página de privacidade com as cláusulas-padrão
    da Resolução CD/ANPD 19/2024. \textbf{Para a IA vão textos sem
    identificação}: a série e o conteúdo, nunca nome, escola ou identificador
    do estudante. Os registros técnicos mascaram CPF, e-mail, nome e textos
    livres; a telemetria usa identificador opaco.
  \item[Limitações] \textbf{Não há prazo nem rotina de retenção}, de exclusão
    ou de anonimização: desligar alguém o desativa, sem apagar (questão aberta
    no Cap. 7 da \cd{PU-DOC-001}, ``decisão de produto e de contrato''). O CPF é
    guardado sem cifra adicional (\adrr{ADR-33}). O avatar da Ajuda não está
    listado entre os suboperadores da página de privacidade. A afirmação de
    que o provedor não treina modelos com os dados não cita a cláusula
    contratual correspondente. Só há diagrama de fluxo de dados do
    check-in de bem-estar; não há um diagrama geral.
  \item[Evidência] \arq{apps/\B{}hub/\B{}app/\B{}privacidade/\B{}page.tsx};
    \arq{apps/\B{}api/\B{}src/\B{}observabilidade/\B{}redaction.ts}; \arq{docs/INFRA.md};
    \cd{PU-DOC-001}, Cap. 7 (região e transferência) e Cap. 11.
  \item[Lacuna] Jurídico/DPO: definir retenção e exclusão; incluir o avatar
    entre os suboperadores; anexar os termos do Google Cloud sobre uso de
    dados (o serviço Vertex AI, por contrato, não usa dados do cliente para
    treinar modelos, mas a cláusula deve ser citada, não presumida).
    Engenharia: desenhar o diagrama geral de fluxo de dados.
\end{campos}
\end{resposta}

\begin{atencao}
\textbf{Ponto sensível para o TR}: os dados ficam fora do Brasil
(\cd{us-east1}) e a IA roda em endpoint global. Se a contratante exigir
residência no país, isso é uma mudança de infraestrutura, e o modelo de IA
usado hoje não está disponível na região de São Paulo. O TR deve dizer se
admite transferência internacional com as salvaguardas da ANPD ou se exige
residência, porque a resposta muda custo e prazo.
\end{atencao}

% ---------------------------------------------------------------------------
\section{T07 · Trilha de auditoria}
\label{sec:t07}

\begin{resposta}{T07}{Eventos auditáveis}{\parcial}
\begin{campos}
  \item[Hoje] Três mecanismos. (1) \textbf{Auditoria de recursos}: ligar ou
    desligar produto e missão, ativar ou fixar versão, alterar quota e
    configuração de IA. Cada registro guarda recurso, escola, ação, versão,
    motivo, valor anterior e novo, autor (com o nome congelado) e data/hora;
    alteração é recusada pelo banco. (2) \textbf{Histórico da evidência
    pedagógica}: cada marca, recálculo, confirmação e ajuste do professor
    gera um evento com valor anterior, valor novo, autor e data; o banco só
    aceita inclusão (alteração recusada, e o serviço não tem permissão de
    apagar). (3) \textbf{Registros técnicos} estruturados com código estável,
    identificador de correlação da requisição, escola e perfil, retidos por 30
    dias. Algumas missões guardam ainda o histórico de comandos do professor
    (quem pausou ou liberou e quando).
  \item[Limitações] \textbf{Não há auditoria de}: login e acesso (só registro
    técnico), criação, edição ou publicação de missão, aprovação de conteúdo,
    ações administrativas sobre usuários e leitura de dados sensíveis. A
    auditoria de recursos permite exclusão no banco. Nenhuma tela ou rota
    consulta o histórico da evidência; ele só é lido direto no banco.
  \item[Planejado] Destino de registros de auditoria por contrato com
    \textbf{400 dias} de retenção, reunindo eventos de autenticação,
    importação, publicação e atividade administrativa da nuvem
    (\adrr{ADR-57}/\adrr{ADR-58}, risco R-05, meta antes de nov./2026). Hoje
    está pendente.
  \item[Evidência] \arq{apps/\B{}api/\B{}src/\B{}recursos/\B{}auditoria.ts}; tabelas
    \cd{recursos\_auditoria} e \cd{evidencias\_eventos} (\cd{PU-DOC-001},
    Cap. 8); \cd{PU-DOC-001}, Cap. 17.
\end{campos}
\end{resposta}

\begin{paraotr}
A solução deve registrar, de forma não alterável pelo usuário e consultável
pela contratante, no mínimo: autenticação, alterações de cadastro e de
vínculo, publicação de atividade, aprovação de conteúdo gerado por IA e
revisão de avaliação, com ator, data/hora, objeto, valor anterior, valor novo,
justificativa (quando houver) e identificador de correlação, retidos por pelo
menos \textit{[prazo a definir]}. Aceite: exportação da trilha de um roteiro
de homologação executado com dados sintéticos.
\end{paraotr}

% ---------------------------------------------------------------------------
\section{T08 · Falha de conexão ou de dispositivo}
\label{sec:t08}

\begin{resposta}{T08}{Progresso, retomada e perda}{\parcial}
\begin{campos}
  \item[Hoje] \textbf{O estado da missão vive no servidor.} Cada jogada que
    vale é uma chamada à API, gravada na hora na sessão do estudante; não se
    espera o fim da missão para salvar. Recarregar a página, fechar o
    navegador ou trocar de computador retoma no passo salvo, comportamento
    registrado nas missões do catálogo (\cd{PU-DOC-001}, Cap. 12). A conclusão
    é gravada uma única vez por estudante e missão: repetir não cria registro
    novo. As rotas do estudante são idempotentes (repetir uma jogada devolve o
    estado atual), e algumas missões mantêm uma fila de reenvio para quedas
    curtas de rede. Os eventos de telemetria têm identificador gerado no
    navegador e são gravados sem duplicar; ficam num \emph{buffer} de até 500
    eventos, enviados em lotes a cada 10 segundos, com novas tentativas
    espaçadas até 60 segundos e envio de emergência ao fechar a aba.
  \item[Limitações] \textbf{Exige conexão contínua}: não há modo
    \emph{offline}, por desenho, porque a correção é feita no servidor e o
    gabarito nunca vai ao navegador. Uma jogada feita sem rede não é aceita
    até reconectar; o reenvio automático só existe nas missões que o
    implementaram. A fila de telemetria fica em memória e se perde se a aba
    for fechada à força ou faltar energia antes do envio; isso afeta só
    telemetria, não o progresso nem as evidências, que já foram gravados na
    jogada.
  \item[Planejado] Guardar a fila de telemetria no armazenamento do navegador
    (IndexedDB), sem prazo.
  \item[Evidência] \arq{packages/\B{}sdk/\B{}src/\B{}telemetry.ts} e
    \arq{telemetry.test.ts}; \arq{apps/\B{}api/\B{}src/\B{}progress/\B{}progresso.controller.ts};
    \cd{PU-DOC-001}, Cap. 10 (onde pode haver perda) e Cap. 12 (retomada por
    missão).
\end{campos}
\end{resposta}

\begin{paraotr}
A solução deve preservar o progresso e as respostas já enviadas pelo
estudante em caso de queda de conexão, fechamento do navegador ou troca de
equipamento, permitindo retomar do último passo registrado, sem duplicar
participação nem evidência. Aceite: roteiro em homologação com interrupção de
rede e reabertura em outro equipamento, por missão.
\end{paraotr}

% ---------------------------------------------------------------------------
\section{T09 · Escala, disponibilidade, monitoramento e suporte}
\label{sec:t09}

\begin{resposta}{T09}{O que pode ser comprovado}{\parcial}
\begin{campos}
  \item[Hoje] \textbf{Monitoramento}: registros estruturados, rastreamento
    distribuído, métricas e erros de navegador estão ativos no ambiente de
    desenvolvimento (\cd{PU-DOC-001}, Cap. 17). Lá estão no ar 4 sondas de
    disponibilidade a partir de 3 regiões, 12 alertas, 2 objetivos de nível de
    serviço (API e hub, 99,5\% sem erro de servidor em 30 dias), fila de
    mensagens mortas da telemetria e painel de operação. Objetivos definidos
    para produção: disponibilidade de hub e API de 99,5\%; 99\% das
    requisições abaixo de 800\,ms (p95); 98\% das chamadas de IA concluídas;
    99\% da telemetria ingerida em menos de 5 minutos. \textbf{Contingência}:
    banco com alta disponibilidade e recuperação a um ponto no tempo; hub e
    API voltam à revisão anterior em segundos; uma missão volta à versão
    anterior sem novo \emph{build}; versão nova pode ser testada numa escola
    antes da rede (Mapa de Produto, H10.3.1 e H10.3.2).
  \item[Limitações] \textbf{Nenhum teste de carga foi executado}. O
    \cd{PU-EST-001} é um plano (versão 0.1, 23/09/2026) com sete testes,
    critérios de aprovação e roteiro, e o risco R-01 o classifica como ``o
    maior risco técnico aberto''. A arquitetura foi dimensionada para cerca de
    10 mil estudantes e 2 a 4 mil simultâneos por contrato. Homologação e
    produção ainda não têm os alertas aplicados. \textbf{Não existe
    compromisso de suporte} (canais, horários, severidades, prazos de resposta
    e de restabelecimento) nem janela de manutenção formal. RPO, RTO e
    retenção de cópias de segurança não estão documentados.
  \item[Planejado] Executar os testes do \cd{PU-EST-001}; aplicar alertas em
    homologação e produção; alerta por consumo do orçamento de erro.
  \item[Evidência] \cd{PU-EST-001}; \cd{PU-DOC-001}, Cap. 6 (reversões), 17
    (SLOs e alertas) e 19 (R-01).
  \item[Lacuna] Comercial/Operação: definir o modelo de suporte. Engenharia:
    executar o teste de carga e declarar RPO/RTO.
\end{campos}
\end{resposta}

\begin{paraotr}
Disponibilidade mensal mínima de 99,5\% para as funções de estudante e
professor, medida por sonda externa e relatada mensalmente; tempo de resposta
p95 abaixo de 800\,ms em carga equivalente ao pico declarado pela contratante
(\textit{[número]} estudantes simultâneos); RPO e RTO declarados pela
contratada e comprovados em exercício de recuperação. Aceite: relatório de
teste de carga com metodologia e resultados antes da produção, e relatório
mensal de disponibilidade e incidentes.
\end{paraotr}

% ---------------------------------------------------------------------------
\section{T10 · Indicadores de adoção}
\label{sec:t10}

\begin{resposta}{T10}{Eventos e painéis de uso}{\parcial}
\begin{campos}
  \item[Hoje] \textbf{Relatório da diretoria} (Mapa de Produto, H07.2.1 e
    H07.2.2, \adrr{ADR-82}): por missão e por escola, ``turmas abertas'',
    ``turmas encerradas'' e ``estudantes com registro'', com exportação em
    planilha; o gestor regional vê a própria regional e o gestor do contrato
    escolhe a regional. \textbf{Cadastro de usuários}: filtro ``ativos /
    inativos / sem primeiro acesso''. \textbf{Análise de produto}: sessões,
    envolvimento e tempo de uso por atividade, só para uma missão
    (Unoovista). \textbf{Uso de IA}: custo e consumo por produto, modelo e
    escola, em vários períodos, mas só para a administração Unooria. Os
    números são lidos na hora da consulta. Eventos coletados: liberação e
    encerramento da missão na turma, sessões e jogadas por estudante,
    conclusão, telemetria das missões, primeiro acesso.
  \item[Limitações] O relatório filtra só por regional: \textbf{não há filtro
    por período, bimestre, ano/série ou turma}. \textbf{Não existe registro de
    último acesso}, nem contagem de professores que acessaram, de missões
    configuradas, iniciadas e concluídas por período, de participação ao
    longo do tempo, nem lista de pendências ou de baixa participação. Os
    períodos letivos existem no cadastro (importados do OneRoster), mas não
    são usados nos relatórios. Cadastrados e ativos se distinguem só no
    filtro de usuários, não num indicador.
  \item[Tratamento] Duplicidade: um registro de progresso por estudante e
    missão, e reimportação idempotente. Transferência: a identidade acompanha
    o estudante, o que ele produziu fica na escola de origem, e a escola nova
    vê o histórico.
  \item[Planejado] H07.3.2 ``Ver os números do dia'' (a fazer): turmas com
    missão aberta, estudantes que entraram e missões concluídas no dia, por
    regional. É o compromisso V1-01 (Capítulo~\ref{cap:v1}).
  \item[Evidência] \arq{apps/\B{}hub/\B{}app/\B{}(plataforma)/\B{}admin/\B{}relatorios/\B{}};
    \arq{apps/\B{}api/\B{}src/\B{}rostering/\B{}escolas.controller.ts}; Mapa de Produto,
    épicos E07 e E03.
\end{campos}
\end{resposta}

% ---------------------------------------------------------------------------
\section{T11 · Pré-visualização integral pelo professor}
\label{sec:t11}

\begin{resposta}{T11}{Preview da missão contextualizada}{\naoexiste}
\begin{campos}
  \item[Hoje] \textbf{Não existe atualmente.} A missão só abre para o perfil
    estudante: professor e gestor são levados de volta à própria área. O
    professor conta com a ficha da missão, o guia da aula e a revisão item a
    item do conteúdo gerado (P04), além de ``zerar a turma'' para reiniciar
    uma rodada. Algumas telas de preparo mostram prévia de textos e
    legendas, mas nenhuma permite jogar a missão.
  \item[Limitações] Como o recurso não existe, não há garantia a descrever
    sobre participação ou evidência geradas no preview.
  \item[Planejado] Não há história no Mapa de Produto nem item no
    \emph{backlog}. É o compromisso V1-02 (Capítulo~\ref{cap:v1}), onde está a
    proposta de implementação.
  \item[Evidência] \arq{apps/\B{}hub/\B{}app/\B{}exp/\B{}[expId]/\B{}page.tsx} (restrição ao
    estudante).
\end{campos}
\end{resposta}

% ===========================================================================
%  CAPITULO 5 — Editorial, Formação e Implantação (E01–E03)
% ===========================================================================
\chapter{Editorial, Formação e Implantação}
\label{cap:editorial}

As três perguntas deste bloco dependem, em grande parte, de material que não
está no repositório da plataforma. O que segue separa o que a plataforma
comprova do que o time editorial e o de implantação precisam anexar.

% ---------------------------------------------------------------------------
\section{E01 · Coleção didática e missões, por ano e bimestre}
\label{sec:e01}

\begin{resposta}{E01}{Matriz curricular integrada}{\foradorepo}
\begin{campos}
  \item[Plataforma] Cada missão é identificada por ano/série e unidade (1 a
    4), na grade de 12 séries $\times$ 4 unidades = 48 posições. Hoje há
    \textbf{14 missões publicadas}, cobrindo as 12 séries na unidade 1, mais
    a unidade 2 do 6º ano e a unidade 4 do 8º ano (Apêndice~\ref{ap:catalogo},
    com as habilidades BNCC de cada régua). Cada missão tem guia de aula com
    objetivos e habilidades BNCC (P01). Novas missões e versões são
    publicadas sem interromper o uso: cada versão tem endereço próprio,
    pode ser testada numa escola antes da rede e revertida (T09); toda
    missão nasce com interruptor de disponibilidade por rede
    (\adrr{ADR-87}).
  \item[Limitações] A plataforma não tem o vínculo com capítulos ou aulas do
    livro (só texto livre em alguns guias) nem a equivalência formal unidade
    = bimestre. A estrutura ``sete aulas com o livro e a oitava com a
    missão'' não aparece no repositório. Três missões não têm régua de
    avaliação por habilidade (Foto Conta, Rastro do Boato, Manifesto da IA).
  \item[Lacuna] Editorial: a matriz ano $\times$ bimestre com tema, objetivos,
    habilidades, capítulos/aulas e missão; a lista de componentes entregues
    (livro do aluno, do professor, recursos digitais); o processo de
    atualização curricular; amostras.
\end{campos}
\end{resposta}

% ---------------------------------------------------------------------------
\section{E02 · Orientações e materiais de apoio ao professor}
\label{sec:e02}

\begin{resposta}{E02}{Onde está cada orientação}{\parcial}
\begin{campos}
  \item[Na plataforma] (1) \textbf{Ficha da missão}: como funciona, preparo,
    competência e formato. (2) \textbf{Guia da aula} de toda missão
    publicada, com versão para impressão: objetivos, habilidades, preparo
    técnico e pedagógico, roteiro minuto a minuto, debate, mediação,
    síntese, rubrica, desdobramentos, cuidados, adaptações e dúvidas (P01).
    (3) \textbf{Exemplos de personalização}: o banco padrão de cada missão.
    (4) \textbf{Interpretação dos resultados}: critérios por nível no painel
    de habilidades. (5) \textbf{Ajuda} com assistente por voz e avatar para
    professores e gestores, cuja base de conhecimento é gerada do Mapa de
    Produto, dos guias e dos manifestos (\adrr{ADR-132}, \adrr{ADR-158}).
    (6) \textbf{Agente de IA do professor} sobre BNCC e o acervo de livros
    (parcial, P08).
  \item[Cobertura] O guia é obrigatório em toda missão publicada (um teste
    recusa publicar sem ele), então cobre as 14 missões.
  \item[Lacuna] Editorial: listar o que está no livro do professor e o que é
    entregue à parte, para o TR não exigir na plataforma o que está só no
    impresso, e vice-versa.
\end{campos}
\end{resposta}

% ---------------------------------------------------------------------------
\section{E03 · Implantação, formação e apoio às redes}
\label{sec:e03}

\begin{resposta}{E03}{Escopo de serviço}{\foradorepo}
\begin{campos}
  \item[Plataforma] Ferramentas que a implantação usa: uma infraestrutura
    isolada por contrato; carga de escolas, turmas e pessoas por OneRoster ou
    planilha, com relatório de erros (T03); criação de escolas, regionais e
    gestores pela cadeia de papéis (T04); interruptor de produto e missão por
    rede, regional ou escola (\adrr{ADR-87}); relatório da diretoria (T10); a
    Ajuda in-app.
  \item[Limitações] Não há no repositório plano de implantação, modelo de
    formação (público, carga horária, modalidade, certificação), canais de
    suporte técnico e pedagógico, regra de escalonamento nem entregáveis de
    fiscalização. Também não há registro de participação em formação na
    plataforma.
  \item[Lacuna] Implantação/Comercial: etapas e marcos de aceite, pré-requisitos
    da Secretaria, modelo de formação, canais e horários de suporte (ver
    T09), relatórios de fiscalização.
\end{campos}
\end{resposta}

\begin{paraotr}
A contratada deve apresentar, no início do contrato, plano de implantação com
etapas, responsáveis e marcos de aceite (carga cadastral validada, acesso de
gestores, formação inicial concluída, primeira missão aplicada por escola), e
entregar relatórios mensais de adoção (V1-01), de formação (participantes e
carga horária) e de atendimento (chamados por severidade e prazo de solução).
\end{paraotr}

% ===========================================================================
%  CAPITULO 6 — Compromissos obrigatórios da V1
% ===========================================================================
\chapter{Compromissos obrigatórios da V1}
\label{cap:v1}

Os seis requisitos abaixo já foram decididos como obrigatórios para a primeira
versão, independentemente do estado atual. Para cada um: status, diferença
para o comportamento esperado, proposta de entrega e plano de aceite.
\textbf{Responsável e data de homologação não foram fixados no repositório}
e ficam marcados como \textit{[a definir]} para o time completar; este
documento não inventa prazo.

\begin{center}
\footnotesize\renewcommand{\arraystretch}{1.35}
\begin{tabularx}{\linewidth}{@{}L{14mm}L{50mm}L{36mm}Y@{}}
\toprule
\textbf{Item} & \textbf{Compromisso} & \textbf{Status} & \textbf{Maior lacuna}\\
\midrule
V1-01 & Indicadores de adoção e uso & \parcial & filtros por período,
  bimestre e turma; último acesso; baixa participação\\
V1-02 & Pré-visualização completa & \naoimpl & o recurso inteiro\\
V1-03 & Recomendações por IA & \naoimpl & o recurso inteiro\\
V1-04 & Recomendação com o tema do setup & \naoimpl & depende de V1-03; o
  tema já está gravado\\
V1-05 & Revisão docente da avaliação & \implementado & teste e2e de ajuste em
  lote\\
V1-06 & Justificativa e auditoria & \parcial & justificativa opcional;
  histórico sem tela\\
\bottomrule
\end{tabularx}
\end{center}

% ---------------------------------------------------------------------------
\section{V1-01 · Indicadores de adoção e uso}

\begin{resposta}{V1-01}{Adoção por rede, escola, turma e bimestre}{\parcial}
\begin{campos}
  \item[Existe] Relatório da diretoria por missão e escola (turmas abertas,
    turmas encerradas, estudantes com registro), recortado por regional e
    exportável; filtro de usuários ``ativos / inativos / sem primeiro
    acesso''; períodos letivos já importados do OneRoster (T03, T10).
  \item[Diferença] Faltam: filtros por \textbf{período e bimestre}, por
    \textbf{ano/série} e por \textbf{turma}; o registro de \textbf{último
    acesso} de professores e estudantes; a contagem de atividades
    configuradas, iniciadas e concluídas; a \textbf{taxa de participação}
    (estudantes que entraram sobre matriculados na turma com missão liberada)
    e a lista das turmas e escolas com participação abaixo de um limiar; a
    série no tempo. A história H07.3.2 (``números do dia'') cobre parte disso
    e está a fazer.
  \item[Proposta] (1) Gravar data e hora do último acesso por usuário a cada
    login. (2) Um painel de adoção para os três níveis de gestão, com os
    filtros acima, usando os dados que já existem: liberações, atividades
    postas no ar, progresso e conclusão por estudante. (3) Denominador
    explícito: matriculados ativos na data, de modo que transferência e
    desligamento não inflem nem zerem a taxa. (4) Indicador ``cadastrados
    $\times$ com primeiro acesso $\times$ ativos no período''.
  \item[Responsável] Produto + Engenharia \textit{[a definir]}.
  \item[Homologação] \textit{[a definir]}.
  \item[Aceite] Rede sintética com escolas de alta e baixa participação; o
    painel deve mostrar os números esperados, por filtro, e apontar as de
    baixa participação. Suíte e2e nova por papel de gestão.
\end{campos}
\end{resposta}

% ---------------------------------------------------------------------------
\section{V1-02 · Pré-visualização completa do professor}

\begin{resposta}{V1-02}{Experimentar a missão antes da turma}{\naoimpl}
\begin{campos}
  \item[Existe] Nada que permita jogar a missão: hoje ela só abre para o
    perfil estudante (T11). Existem a ficha, o guia e a revisão item a item
    dos conteúdos gerados (P01, P04).
  \item[Diferença] O recurso inteiro. Pontos que a implementação precisa
    resolver: percorrer todas as etapas e ramificações com o conteúdo
    aprovado; permitir avançar as etapas que, na aula, dependem de comando do
    professor; voltar ao preparo para corrigir um item; e não gerar
    participação, conclusão nem evidência.
  \item[Proposta] Um modo \emph{Experimentar como estudante}, aberto a partir
    da atividade, que roda a mesma missão com uma sessão marcada como
    prévia, pertencente ao professor e fora da turma. As rotas de estudante
    reconhecem a marca e não gravam progresso, conclusão, telemetria
    pedagógica nem evidência (ou gravam num espaço separado que nenhum painel
    lê). As etapas comandadas pelo professor ficam disponíveis num controle
    da própria prévia. Um botão leva de volta ao preparo da atividade.
  \item[Limitações] As 14 missões têm rotas próprias de estudante: o custo é
    por missão, e a entrega pode ser escalonada pelas missões do primeiro
    bimestre.
  \item[Responsável] Produto + Engenharia \textit{[a definir]}.
  \item[Homologação] \textit{[a definir]}.
  \item[Aceite] Por missão: o professor percorre a prévia até o fim; depois,
    painel da aula, painel de habilidades, relatórios e exportações da turma
    continuam idênticos ao estado anterior (teste automatizado que compara os
    números antes e depois).
\end{campos}
\end{resposta}

% ---------------------------------------------------------------------------
\section{V1-03 · Recomendações pedagógicas por IA}

\begin{resposta}{V1-03}{Sugestões de intervenção}{\naoimpl}
\begin{campos}
  \item[Existe] As entradas já existem: resultado por habilidade na turma
    (atingiu, em processo, não atingiu, não avaliado), os critérios por nível
    de cada habilidade, o guia da aula (objetivos, rubricas, desdobramentos)
    e a infraestrutura de IA com quota, guardrails e medição de custo. Não há
    geração de recomendação (P08).
  \item[Diferença] O recurso inteiro e o processo pedagógico em volta dele:
    autoria e versionamento dos prompts pela Pedagogia, conjunto de
    avaliação com casos aprovados por pedagogos, e regra de apresentação que
    separe dado observado de sugestão.
  \item[Proposta] Gerar, sob demanda do professor no painel de habilidades,
    uma recomendação em formato fechado: \emph{o que os dados mostram}
    (somente contagens e critérios, copiados do painel, nunca escritos pela
    IA) e \emph{o que sugerimos} (intervenções ligadas a cada habilidade em
    dificuldade, citando o código). Prompt e diretrizes ficam no servidor,
    com versão gravada junto de cada recomendação. Nenhum dado nominal de
    estudante vai ao modelo. Consome a quota do professor.
  \item[Responsável] Pedagogia (prompts e casos de avaliação) + IA
    \textit{[a definir]}.
  \item[Homologação] \textit{[a definir]}.
  \item[Aceite] Conjunto de casos de turma sintética revisado por pedagogos,
    com critérios de aprovação (pertinência, ausência de afirmação sem
    suporte nos dados, adequação à série) e taxa mínima acordada; a
    recomendação nunca afirma número que não esteja no painel.
\end{campos}
\end{resposta}

% ---------------------------------------------------------------------------
\section{V1-04 · Recomendações considerando a contextualização}

\begin{resposta}{V1-04}{Tema do setup na recomendação}{\naoimpl}
\begin{campos}
  \item[Existe] O tema, os tópicos e o componente curricular informados pelo
    professor ao criar a atividade já ficam gravados nela (P02), assim como
    a atividade de onde veio cada resultado.
  \item[Diferença] Depende de V1-03. Falta levar esses campos à geração da
    recomendação, por atividade.
  \item[Proposta] A recomendação de V1-03 recebe, além dos resultados, o
    componente, o tema e os tópicos da atividade, e propõe intervenções no
    contexto escolhido (no exemplo, desinformação aplicada a ``vacinas'' em
    Ciências), sem alterar habilidade nem critério.
  \item[Responsável] Pedagogia + IA \textit{[a definir]}.
  \item[Homologação] \textit{[a definir]}, junto de V1-03.
  \item[Aceite] Para a mesma turma sintética com dois temas diferentes, as
    recomendações mudam de contexto e mantêm as mesmas habilidades e os
    mesmos números.
\end{campos}
\end{resposta}

% ---------------------------------------------------------------------------
\section{V1-05 · Revisão docente da avaliação}

\begin{resposta}{V1-05}{O professor tem a decisão final}{\implementado}
\begin{campos}
  \item[Existe] O professor que leciona na turma ajusta o nível de qualquer
    habilidade de um ou de vários estudantes, por atividade e rodada, a
    qualquer momento, inclusive com a missão encerrada, e pode marcar ``não
    participou''. O ajuste passa a ser a origem ``professor'' e nada o
    sobrescreve depois. Também confirma, para a turma, níveis projetados.
    Gestores só leem.
  \item[Diferença] O Mapa de Produto marca a história como parcial porque
    falta o teste e2e do ajuste em lote e do histórico (H14.4.2). Ressalva:
    \emph{zerar a turma} apaga também o ajuste vigente (\adrr{ADR-155});
    o evento do ajuste continua no histórico (V1-06).
  \item[Responsável] Produto + Pedagogia \textit{[a definir]}.
  \item[Homologação] Pronto para demonstração; falta o teste citado.
  \item[Aceite] Em homologação: o professor ajusta um nível; o painel de
    habilidades, o perfil do estudante e os relatórios passam a mostrar o
    nível do professor.
  \item[Evidência] \arq{apps/\B{}api/\B{}src/\B{}evidencias/\B{}evidencias-professor.controller.ts};
    \arq{apps/\B{}hub/\B{}components/\B{}evidencias/\B{}ajustar-selecionados.tsx}.
\end{campos}
\end{resposta}

% ---------------------------------------------------------------------------
\section{V1-06 · Justificativa e auditoria obrigatórias}

\begin{resposta}{V1-06}{Trilha da revisão}{\parcial}
\begin{campos}
  \item[Existe] Cada ajuste gera um evento permanente com o valor original,
    o valor novo, o autor e a data e hora; o banco só aceita inclusão (a
    alteração é recusada e o serviço não tem permissão de apagar), e o
    histórico sobrevive a \emph{zerar a turma}. Há um motivo por lista
    (presente sem participar, problema técnico, legenda revisada, outro).
  \item[Diferença] (1) O motivo é \textbf{opcional} (a tela oferece ``sem
    motivo'') e não há texto livre de justificativa. (2) \textbf{Nenhuma tela
    ou rota consulta o histórico}: hoje ele só é lido no banco. (3) Não há
    exportação da trilha.
  \item[Proposta] Tornar o motivo obrigatório, com texto livre obrigatório
    quando ``outro''; mostrar, por estudante e habilidade, o histórico
    completo (original, cada alteração, autor, data e justificativa) ao
    professor e à coordenação; exportação da trilha para a gestão.
  \item[Responsável] Produto + Engenharia/Segurança \textit{[a definir]}.
  \item[Homologação] \textit{[a definir]}.
  \item[Aceite] Ajuste sem justificativa é recusado; após três ajustes
    sucessivos, o histórico mostra o resultado original e as três alterações
    com autor, data e motivo; tentativa de alterar ou apagar o histórico
    direto no banco com o usuário da aplicação é recusada.
\end{campos}
\end{resposta}

% ===========================================================================
%  CAPITULO 7 — Consolidação para o Termo de Referência
% ===========================================================================
\chapter{Consolidação para o Termo de Referência}
\label{cap:consolidacao}

\section{O que o TR pode afirmar com segurança}

Itens com evidência no código e em testes, que podem ser exigidos como
requisito sem risco de a solução atual não atender:

\begin{itemize}
  \item Acesso 100\% web, sem instalação, com interface responsiva.
  \item Carga cadastral no padrão OneRoster 1.2 (CSV), por interface e API,
    com importação parcial, relatório de erros e carga idempotente.
  \item Seis perfis com autorização aplicada no servidor e isolamento lógico
    por escola; isolamento físico por contrato.
  \item Ficha e guia de aula por missão, com objetivos, habilidades BNCC,
    roteiro e rubrica.
  \item Personalização por IA com aprovação obrigatória do professor antes do
    uso, sem alterar a estrutura pedagógica nem os critérios de avaliação.
  \item Avaliação por habilidade com regras no servidor, histórico de rodadas
    e revisão pelo professor, que prevalece sobre a classificação
    automática.
  \item Proteção do texto do estudante: nada de identificação enviada à IA;
    textos livres fora de painéis e registros técnicos.
  \item Progresso salvo no servidor a cada interação, com retomada.
\end{itemize}

\section{O que o TR não deve afirmar ainda}

\begin{limitacao}
\begin{itemize}
  \item Ganhos de aprendizagem, engajamento ou aceitação (P10 sem relatório).
  \item Capacidade comprovada de carga (nenhum teste executado, T09).
  \item Conformidade declarada com WCAG AA (falta a avaliação manual, P09).
  \item Login federado por OIDC ou SAML em operação (T02).
  \item Duplicação ou compartilhamento de missões personalizadas (P03).
  \item Anexar documentos para orientar a IA (P02).
  \item Funcionamento \emph{offline} (T08).
  \item Recomendações pedagógicas por IA e pré-visualização (V1-02 a V1-04:
    são compromissos, não funções existentes).
\end{itemize}
\end{limitacao}

\section{Decisões que o TR precisa tomar}

\begin{decisao}
\begin{enumerate}
  \item \textbf{Residência de dados}: admitir transferência internacional com
    as salvaguardas da ANPD ou exigir dados no Brasil (T06). A segunda opção
    muda infraestrutura, custo e o modelo de IA disponível.
  \item \textbf{Retenção e exclusão}: prazos por categoria de dado e o
    destino ao fim do contrato (T06).
  \item \textbf{Revogação de acesso}: prazo máximo entre desligar uma pessoa e
    ela perder acesso, inclusive em sessão aberta (T02).
  \item \textbf{Integração cadastral}: se basta a rede enviar as cargas ou se
    será exigida sincronização automática (T03).
  \item \textbf{Nível de serviço e suporte}: disponibilidade, horários,
    severidades e prazos (T09, E03).
  \item \textbf{Equipamento de referência}: navegadores, telas e banda mínimos
    a declarar no edital (T01).
\end{enumerate}
\end{decisao}

\section{Pendências por time}

\begin{center}
\footnotesize\renewcommand{\arraystretch}{1.35}
\begin{tabularx}{\linewidth}{@{}L{34mm}X@{}}
\toprule
\textbf{Time} & \textbf{O que precisa entregar para fechar o TR}\\
\midrule
\endhead
Produto / Pedagogia & Relatório da PoC (P10); capturas de ficha e guia (P01);
  par entrada/saída de personalização (P02); prompts e casos de avaliação
  das recomendações (V1-03, V1-04); responsáveis e datas da V1.\\
Editorial & Matriz ano $\times$ bimestre com livro e missão; componentes
  entregues; o que está no livro e o que está na plataforma (E01, E02).\\
Engenharia & Matriz de compatibilidade (T01); teste de carga do
  \cd{PU-EST-001} (T09); RPO e RTO; diagrama geral de fluxo de dados (T06);
  relatório atualizado de acessibilidade e avaliação manual (P09);
  configuração dos filtros do provedor de IA (T05); auditoria de
  autenticação e publicação (T07); V1-01, V1-02 e V1-06.\\
Jurídico / DPO & Retenção e exclusão; suboperadores (incluir o avatar da
  Ajuda); cláusula do provedor de IA sobre uso de dados (T06).\\
Implantação / Comercial & Plano de implantação, formação e suporte (E03,
  T09).\\
\bottomrule
\end{tabularx}
\end{center}

\section{Recomendações de forma para o TR}

\begin{itemize}
  \item Escrever cada requisito como resultado observável com critério de
    aceite, como nas caixas ``Proposta de redação'' deste documento, sem
    nomear tecnologia, nuvem ou marca.
  \item Exigir demonstração em homologação com dados sintéticos para todo
    requisito declarado como atendido.
  \item Tratar acessibilidade e proteção de dados como exigência legal, não
    como item de \emph{roadmap}.
  \item Manter fora do escopo mínimo as ideias que não são compromisso da V1
    (duplicação e compartilhamento, anexos, modo \emph{offline}, SAML) ou
    marcá-las como opcionais pontuáveis.
\end{itemize}


\appendix
\renewcommand{\thechapter}{\Alph{chapter}}
% ===========================================================================
%  APENDICE A — Catálogo de missões em 30/10/2026
% ===========================================================================
\chapter{Catálogo de missões em 30/10/2026}
\label{ap:catalogo}

Esta tabela foi lida do manifesto público de cada missão, o arquivo
\arq{public/\B{}experiencia.json} da pasta dela em \arq{apps/\B{}experiencias/}.
Os manifestos não declaram habilidades BNCC: os códigos abaixo vêm das réguas
de avaliação, em \arq{apps/\B{}api/\B{}src/\B{}evidencias/\B{}fichas/}. Todas
pertencem ao produto \emph{Entre Redes}. ``Unid.'' é a unidade do ano (1 a 4).

\begin{center}
\footnotesize\renewcommand{\arraystretch}{1.3}
\begin{tabularx}{\linewidth}{@{}L{20mm}L{52mm}C{10mm}C{9mm}L{13mm}X@{}}
\caption{Missões publicadas}\label{tab:catalogo}\\
\toprule
\textbf{Código} & \textbf{Missão} & \textbf{Série} & \textbf{Unid.} &
  \textbf{Versão} & \textbf{Habilidades na régua}\\
\midrule
\endfirsthead
\toprule
\textbf{Código} & \textbf{Missão} & \textbf{Série} & \textbf{Unid.} &
  \textbf{Versão} & \textbf{Habilidades na régua}\\
\midrule
\endhead
\cd{exp01ef1b} & O Que a Foto Conta & 1º EF & 1 & 1.0.11 & sem régua\\
\cd{exp02ef1b} & O Enigma da Capa da Rádio & 2º EF & 1 & 2.2.3 & EF02CO01, EF02CO05\\
\cd{exp03ef1b} & Nina e a Notícia que Desapareceu & 3º EF & 1 & 0.2.0 & EF03CO01 a 03, EF15CO08\\
\cd{exp04ef1b} & Jornal da Escola: A Edição de Bento & 4º EF & 1 & 3.0.4 & EF04CO08, EF15CO02, EF15CO04\\
\cd{exp05ef1b} & Quem Mandou Isso? & 5º EF & 1 & 0.6.4 & EF05CO03, EF05CO04, EF05CO08\\
\cd{exp06ef1b} & O Júri Estadual de Projetos & 6º EF & 1 & 0.8.4 & EF06CO01, EF06CO03\\
\cd{exp06ef2b} & O Rastro do Boato & 6º EF & 2 & 0.1.11 & sem régua\\
\cd{exp07ef1b} & A Seleção & 7º EF & 1 & 1.3.4 & EF07CO08, EF07CO11\\
\cd{exp08ef1b} & Unoovista: O Conteúdo que não Chegou & 8º EF & 1 & 1.5.5 & EF08CO07, EF08CO09, EF08CO11\\
\cd{exp08ef4b} & O Manifesto da IA & 8º EF & 4 & 0.1.11 & sem régua\\
\cd{exp09ef1b} & O Dossiê Invisível & 9º EF & 1 & 0.1.0 & EF09CO07, EF09CO08\\
\cd{exp01em1b} & O Assistente Invisível & 1ª EM & 1 & 1.0.6 & EM13CO01, 03, 18, 20\\
\cd{exp02em1b} & O Efeito Amplificador & 2ª EM & 1 & 2.3.3 & EM13CO08, 14, 24, 26\\
\cd{exp03em1b} & O Porta-Voz dos Dados & 3ª EM & 1 & 1.0.1 & EM13CO12, 13, 14, 21, 22, 26\\
\bottomrule
\end{tabularx}
\end{center}

\begin{nota}
Fora do catálogo: a versão anterior da missão da 2ª série do EM (Missão
Tempestade, \cd{exp02em1b-v0}), mantida só para as turmas que já a tinham
liberado. A plataforma tem ainda ferramentas do professor que não são missões
(plano de aula, plano de ensino pela BNCC, check-in de bem-estar) e um jogo 3D
(Vila do Lume). Cobertura atual: 14 de 48 posições da grade de 12 séries
$\times$ 4 unidades.
\end{nota}


\end{document}
